%%%%%%%%%%%%%%%%%%%%%%%%%%%%%%%%%%%%%%%%%%%%%%%%%%%%%%%%%%%%%%%%%%%%%%%%%%%%%%%%%%%%%%%%%%%%%%%%%
%%
%%
%%
%% Please do minimal aesthetic changes to the TeX file. 
%% All the changes in the TeXmac file (quasi-linear-final-candidate.tm)
%%
%%
%% Arxiv version : 2017-02-10
%%%%%%%%%%%%%%%%%%%%%%%%%%%%%%%%%%%%%%%%%%%%%%%%%%%%%%%%%%%%%%%%%%%%%%%%%%%%%%%%%%%%%%%%%%%%%%%%%%%
%%%%%%%%%%%%%%%%%%%%%%%%%%%%%%%%%%%%%%%%%%%%%%%%%%%%%%%%%%%%%%%%%%%%%%%%%%%%%%%%%%%%%%%%%%%%%%%%%


%\documentclass{amsart}
%\usepackage[english]{babel}
%\usepackage{amssymb,stmaryrd,enumerate,mathrsfs,bbm,theorem}

\documentclass[11pt,reqno]{amsart}
\usepackage{geometry}
\usepackage{hyperref}
\usepackage[english]{babel}

\geometry{a4paper, margin=1in}

\usepackage{amssymb}
\usepackage{stmaryrd}
\usepackage{enumerate}
\usepackage{mathrsfs}
\usepackage{bbm}
%\usepackage{theorem}
\usepackage{amsthm}

%%%%%%%%%% Start TeXmacs macros
\catcode`\<=\active \def<{
\fontencoding{T1}\selectfont\symbol{60}\fontencoding{\encodingdefault}}
\catcode`\>=\active \def>{
\fontencoding{T1}\selectfont\symbol{62}\fontencoding{\encodingdefault}}
\newcommand{\assign}{:=}
\newcommand{\cdummy}{\cdot}
\newcommand{\mathD}{\mathrm{D}}
\newcommand{\mathd}{\mathrm{d}}
\newcommand{\nocomma}{}
\newcommand{\nospace}{}
\newcommand{\nosymbol}{}
\newcommand{\precprec}{\prec\!\!\!\prec}
\newcommand{\tmdummy}{$\mbox{}$}
\newcommand{\tmop}[1]{\ensuremath{\operatorname{#1}}}
\newcommand{\tmtextit}[1]{{\itshape{#1}}}
\newenvironment{enumeratenumeric}{\begin{enumerate}[1.] }{\end{enumerate}}
\newenvironment{itemizedot}{\begin{itemize} \renewcommand{\labelitemi}{$\bullet$}\renewcommand{\labelitemii}{$\bullet$}\renewcommand{\labelitemiii}{$\bullet$}\renewcommand{\labelitemiv}{$\bullet$}}{\end{itemize}}

%\newtheorem{definition}{Definition}
%\newtheorem{lemma}{Lemma}
%{\theorembodyfont{\rmfamily}\newtheorem{remark}{Remark}}
%\newtheorem{theorem}{Theorem}
\newtheorem{theorem}{Theorem}
\newtheorem{definition}[theorem]{Definition}
\newtheorem{lemma}[theorem]{Lemma}
\newtheorem{remark}[theorem]{Remark}


%%%%%%%%%% End TeXmacs macros

\newcommand{\CC}{\mathscr{C} \hspace{.1em}}
\newcommand{\CD}{\mathscr{D} \hspace{.1em}}
\newcommand{\CF}{\mathscr{F} \hspace{.1em}}
\newcommand{\CB}{\mathscr{B} \hspace{.1em}}
\newcommand{\CA}{\mathscr{A} \hspace{.1em}}
\newcommand{\DD}{\mathscr{D} \hspace{.1em}}
\newcommand{\LL}{\mathscr{L} \hspace{.2em}}
\newcommand{\SSS}{S{\hspace{.1em}}}
\newcommand{\PP}{\mathscr{P} \hspace{.1em}}
\newcommand{\DF}{\mathfrak{D}}
%

%

%

%

%


\begin{document}

\title{Paracontrolled quasilinear SPDEs}

\author{M.~Furlan}
\address{CEREMADE\\
Universit{\'e} Paris Dauphine, France}
\email{furlan@ceremade.dauphine.fr}

\author{ M.~Gubinelli}
\address{IAM \& HCM\\
Universit{\"a}t Bonn, Germany}
\email{gubinelli@iam.uni-bonn.de}

\begin{abstract}
  We introduce a non-linear paracontrolled calculus and use it to renormalise
  a class of singular SPDEs including certain quasilinear variants of the
  periodic two dimensional parabolic Anderson model.
\end{abstract}

{\maketitle}

\section{Introduction}

We show how to renormalise a class of general quasilinear equations of which
one of the simplest examples is the following parabolic SPDE:
\begin{equation}
  \partial_t u (t, x) - a (u (t, x)) \Delta u (t, x) = \xi (x), \qquad u (0,
  x) = u_0 (x), \quad x \in \mathbbm{T}^2, t \geqslant 0,
  \label{eq:ur-formulation}
\end{equation}
with $a : \mathbbm{R} \rightarrow [\lambda, 1]$ for~$\lambda > 0$ a uniformly
bounded $C^3$ diffusion matrix, and $\| a^{(k)} \|_{L^{\infty}} \leqslant 1$
for $k = 0, \ldots, 3$. We assume that $\xi \in \CC \nospace^{\alpha - 2}
(\mathbbm{T}^2)$ with $2 / 3 < \alpha < 1$ where $\CC^{\alpha}
(\mathbbm{T}^2)$ is the Besov space $B^{\alpha}_{\infty, \infty}
(\mathbbm{T}^2)$. This would apply to the space white noise on
$\mathbbm{T}^2$, for example. In this case we only expect that $u (t, \cdot)
\in \CC^{\alpha} (\mathbbm{T}^2)$ and the term $a (u (t, \cdot)) \Delta u (t,
\cdot)$ is not well defined when $2 \alpha - 2 < 0$.
Eq.~(\ref{eq:ur-formulation}) is a quasilinear generalisation of the
two--dimensional periodic parabolic Anderson model (PAM).

Let us remark from the start that the framework we will consider below allows
to deal with a class of equations of the form
\begin{equation}
  a_1 (u (t, x)) \partial_t u (t, x) - a_2 (u (t, x)) \Delta u (t, x) = \xi
  (a_3 (u (t, x)), t, x), \qquad x \in \mathbbm{T}^2, t \geqslant 0,
  \label{eq:general-model}
\end{equation}
where $a_1, a_2$ are sufficiently smooth non-degenerate coefficients and $\xi
(z, t, x)$ is a Gaussian process with covariance
\[ \mathbbm{E} [\xi (z, t, x) \xi (z', t', x')] = F (z, z') Q (t - t', x -
   x'), \qquad x, x' \in \mathbbm{T}^2, t, t', z, z' \in \mathbbm{R}, \]
with $F$ a smooth function and $Q$ a distribution of parabolic regularity
$\rho > - 4 / 3$. This includes the space white noise discussed before or a
time white noise with a regular dependence on the space variable or some noise
mildly irregular in space and time.

Also the scalar character of the equation or of the non-linear diffusion
coefficient will not play any specific role and we could consider
vector--valued equations with general diffusion coefficients provided the
template problem~(\ref{eq:template}) below remains uniformly parabolic.

For the sake of clarity and simplicity we will discuss mainly the basic
example~(\ref{eq:ur-formulation}) since this contains already most of the
technical difficulties. The fact that one can handle models as general
as~(\ref{eq:general-model}) can be considered a direct byproduct of the
techniques we will introduce below.

Recently Otto and Weber~{\cite{otto_quasilinear_2016}} and Bailleul, Debussche
and Hofmanova~{\cite{hofmanova}} investigated quasilinear SPDEs in the context
of pathwise methods and in a range of regularities compatible with the ones we
will consider in this paper.
\begin{itemizedot}
  \item In~{\cite{otto_quasilinear_2016}} the authors obtained a priori
  estimates for equations of the form
  \[ \partial_t u (t, x) - a (t, x) \partial_x^2 u (t, x) = f (u (t, x)) \xi
     (t, x), \qquad (t, x) \in \mathbbm{T}^2 \]
  where both space and time variables take values in a one dimensional
  periodic domain and their noise can be white in time but colored in space,
  essentially behaving like a distribution of parabolic regularity in $(- 4 /
  3, 1)$. In order to do so they introduce a specific notion of
  \tmtextit{modelled} function and related estimates. \ \ \ \
  
  \item Bailleul, Debussche and Hofmanov{\'a} in~{\cite{hofmanova}} obtain
  local well-posedness for the generalised parabolic Anderson model equation
  \begin{equation}
    \partial_t u (t, x) - a (u (t, x)) \Delta u (t, x) = g (u (t, x)) \xi (x)
    \label{eq:deb-hof} \qquad t \geqslant 0, x \in \mathbbm{T}^2 .
  \end{equation}
  The authors obtain the same result as the one presented in
  Section~\ref{sec:gen-nonlinear} of our work, without the machinery of
  nonlinear paraproducts introduced here, but using only the basic tools of
  paracontrolled analysis and some clever transformations.
  
  On the other hand, we remark that the apparently innocuous vectorial
  formulation of~(\ref{eq:deb-hof})
  \[ \partial_t u (t, x) - a_{i j} (u (t, x)) \frac{\partial^2}{\partial x_i
     \partial x_j} u (t, x) = g (u (t, x)) \xi \qquad t \geqslant 0, x \in
     \mathbbm{T}^2 \]
  is out of reach of the techniques used in~{\cite{hofmanova}}, while can be
  treated flawlessly in our framework.
\end{itemizedot}


Let us state one simple result which can be obtained via the theory developed
in this paper:

\begin{theorem}
  Fix $2 / 3 < \alpha < 1$. Let $\xi \in \CC^{\alpha - 2} (\mathbbm{T}^2)$ be
  a space white noise with zero average on the torus, and $u_0 \in
  \CC^{\alpha}$ an initial condition. Let $(\xi_{\varepsilon}, u_{0,
  \varepsilon})_{\varepsilon > 0}$ be a family of smooth approximations to
  $\xi, u_0$ obtained by convolution with a rescaled smoothing kernel and
  $u_{\varepsilon}$ the classical solution to the Cauchy problem
  \begin{equation}
    \partial_t u_{\varepsilon} - a (u_{\varepsilon}) \Delta u_{\varepsilon} =
    \xi_{\varepsilon} + \sigma_{\varepsilon} \frac{a' (u_{\varepsilon})}{a
    (u_{\varepsilon})^2}, \qquad u (0) = u_{0, \varepsilon} . \qquad
    \label{eq:ur-formulation-approx}
  \end{equation}
  Then we can choose the constants $(\sigma_{\varepsilon})_{\varepsilon > 0}$
  and a random time $T > 0$ in such a way that the family of r.v.
  $(u_{\varepsilon})_{\varepsilon > 0} \subseteq \LL_T^{\alpha}
  (\mathbbm{T}^2)$ almost surely converge as $\varepsilon \rightarrow 0$ to a
  random element $u \in \LL_T^{\alpha} (\mathbbm{T}^2)$, where
  $\LL_T^{\alpha}$ is the parabolic space $C \left( [0, T], \CC^{\alpha}
  (\mathbbm{T}^2) \right) \cap C^{\alpha / 2} \left( [0, T], \CC^0
  (\mathbbm{T}^2) \right)$.
  
  This element can be characterised as the solution to a paracontrolled
  singular SPDE (see below for more details).
\end{theorem}

In order to devise a suitable formulation of eq.~(\ref{eq:ur-formulation}) and
obtain a theory with $u \in \CC^{\alpha}$ we decompose the non-linear
diffusion term in the l.h.s. with help of Bony's
paraproduct~{\cite{meyer_remarques_1981}} and write
\begin{equation}
  \partial_t u - a (u) \prec \Delta u = \xi + \Phi (u)
  \label{e:paracontrolled-eq}
\end{equation}
with
\begin{equation}
  \Phi (u) \assign a (u) \circ \Delta u + a (u) \succ \Delta u
  \label{e:def-Phi}
\end{equation}
where $\prec, \succ$ are standard paraproducts and $\circ$ denotes the
resonant product (see below for precise definitions). Now the l.h.s. is always
well defined irrespective of the regularity of the function $u$ and the
problem becomes that of controlling the resonant product $a (u) \circ \Delta
u$ appearing in the r.h.s.~. A key point of the analysis put forward below is
that this term can be expected to be of regularity $2 \alpha - 2 > \alpha - 2$
so better than the leading term $\xi$.

Our approach can be described as follows. For an equation of the form
\[ \partial_t u - a_1 (u) \Delta u = a_2 (u) \xi, \]
we consider at first a \tmtextit{parametric} ``template'' problem with
constant coefficients
\begin{equation}
  \partial_t \vartheta (\eta, t, x) - \eta_1 \Delta \vartheta (\eta, t, x) =
  \eta_2 \xi (t, x), \label{eq:template}
\end{equation}
where now $\eta = (\eta_1, \eta_2)$ are fixed numbers. A \tmtextit{nonlinear}
paraproduct $\Pi_{\precprec}$ will allows us to modulate the parametric
solution $\vartheta$ with the coefficient $a (u) = (a_1 (u), a_2 (u))$ as to
capture the most irregular part of the solution $u$ itself. As a consequence,
the paracontrolled Ansatz
\[ u = \Pi_{\precprec} (a (u), \vartheta) + u^{\sharp} \]
will define a regular remainder term $u^{\sharp}$ which solves a standard PDE.
With this decomposition the resonant products appearing in the equation can be
estimated along the lines of the standard paracontrolled arguments introduced
in~{\cite{gubinelli_paracontrolled_2015}} and all the arguments introduced
there can be extended in a straightforward manner to the quasilinear setting.

This approach has been inspired by the parametric controlled Ansatz of Otto
and Weber~{\cite{otto_quasilinear_2016}}. At variance with their approach we
use the parametric Ansatz in the context of the paradifferential calculus and
consider more general noise terms.

Usefulness of paraproducts in the analysis of non-linear PDEs is by now well
established: see for example the seminal paper of
Meyer~{\cite{meyer_remarques_1981}}, the early review of
Bony~{\cite{bony_microlocal_1991}}, the recent books of Alinhac and
G{\'e}rard~{\cite{alinhac_operateurs_1991}} and Bahouri, Chemin, and
Danchin~{\cite{bahouri_fourier_2011}}. Let us mention also the interesting
paper of H{\"o}rmander~{\cite{hormander_nash_moser_1990}} where
paradifferential operators allows to bypass the Nash--Moser fixpoint theorem
in some applications \ where the loss of regularity prevents straightforward
use of standard Banach fixpoint theorem. The main observation in that paper is
that, with the aid of paradifferential operators, it is possible to identify a
``corrected'' problem for which standard Banach fixpoint applies.

Paracontrolled calculus for singular SPDEs has beed introduced by Gubinelli,
Imkeller and Perkowski~{\cite{gubinelli_paracontrolled_2015}} (see also the
lecture notes~{\cite{gubinelli_lectures_2015}}) and used to study various
equations like the KPZ equation~{\cite{gubinelli_kpz_2015}}, the dynamic
$\Phi^4_3$ model~{\cite{catellier_paracontrolled_2013}} in $d = 3$ and its
global well--posedness~{\cite{mourrat_global_2016}}, the spectrum of the
continuous Anderson Hamiltonian in $d = 2$~{\cite{allez_continuous_2015}}. By
using heat--semigroup techniques paracontrolled calculus has been extended to
the manifold context by Bailleul and Bernicot~{\cite{bailleul_heat_2015}}.

Non--linear generalisation of the classic bilinear paraproducts already
appeared in the notion of \tmtextit{paracomposition} introduced by
Alinhac~{\cite{bony_microlocal_1991}}. Non--linear versions of rough paths
have been considered by one of the authors in order to study the Korteweg--de
Vries equation~{\cite{gubinelli_rough_2012}}. Non--linear Young integrals were
used by one of the authors in joint work with
Catellier~{\cite{catellier_averaging_2016}} to study the the regularising
properties of sample paths of stochastic processes processes. See also the
related work of Hu and Le~{\cite{hu_nonlinear_2014}} on differential systems
in H{\"o}lder media. Relevant to this discussion of non-linear variants of
rough paths is the work of Bailleul on rough
flows~{\cite{bailleul_flows_2015}} and their application to
homogeneisation~{\cite{bailleul_rough_2016}}. By looking at the composition $f
(g (x))$ as the action of the distribution $\delta_{g (x)}$ on the function
$f$, non--linear constructions can be linearised at the price of working in
infinite--dimensional spaces: this is the approach chosen by Kelly and
Melbourne to avoid non--linear generalisations of rough path theory in their
study of homogeneisation of fast--slow system with random initial
conditions~{\cite{kelly_deterministic_2014}}. It is worth mentioning also
Kunita's theory of semimartingale vector
fields~{\cite{kunita_stochastic_1984}} which occupy a place in stochastic
analysis quite similar to that which these recent developments occupy in the
landscape of rough paths/paracontrolled distributions theories.

Paracontrolled calculus is currently limited to ``first order'' computations.
This limitation is also ubiquitous in the present work. Even if, in practice,
this is not a big issue, and the calculus is still able to deal with a large
class of problems, it makes the paracontrolled approach less appealing for a
general theory of singular SPDEs. Let us remark that recently Bailleul and
Bernicot~{\cite{bailleul_higher_2016}} developed an higher order version of
the paracontrolled calculus. However, apart from these recent development,
whose impact is still to be assessed, the most general theory for singular
SPDEs has been developed by
Hairer~{\cite{hairer_theory_2014,hairer_singular_2014,friz_course_2014}} under
the name of \tmtextit{regularity structures} theory. Regularity structures are
a vast generalisation of Lyons' rough
paths~{\cite{lyons_differential_1998,lyons_system_2002,lyons_differential_2007}}
which give effective tools to describe non-linear operations acting on certain
spaces of distributions, their renormalization by subtraction of local
singularities and their use to solve singular SPDE. Regularity structures have
been successfully applied to all the models mentioned so
far~{\cite{hairer_theory_2014,hairer_solving_2013}}, to other models like the
Sine--Gordon model~{\cite{hairer_dynamical_2014}} (which however can also be
handled via paracontrolled techniques) and to study weak universality
conjectures~{\cite{hairer_class_2015,hairer_large_2016}}. In their current
instantiation it does not seem possible to solve quasilinear SPDEs via
regularity structures. The results of the present paper hint to the fact that
a non-linear version of regularity structures is conceivable, at least in
principle. Indeed one can imagine models depending on additional parameters
and modelled distributions acting as evaluations of the parametric models at
certain space--time dependent values of the parameters. It would be
interesting to pursue further this intuition.

\

The structure of the paper is the following. In Section~\ref{s:nonlinear} we
introduce our basic tools: the non-linear paraproduct decomposition and some
related commutation lemmas. In Section~\ref{s:paracontrolled} we introduce the
paracontrolled Ansatz which allows to transform the singular
problem~(\ref{eq:ur-formulation}) into a well--behaved PDE. In
Section~\ref{s:fixpoint} we discuss the apriori estimates, the uniqueness of
the solution of the transformed PDE and its continuity w.r.t. the random data
and the initial condition, we introduce also the algebraic structure which
allows to renormalise the model. Section~\ref{s:renorm} deals with the
renormalization of the stochastic data and the construction of the enhanced
noise associated to white noise. Section~\ref{sec:gen-nonlinear} deals with
the extension of the results to more general equations, in particular with
equation~(\ref{eq:deb-hof}) or with noise whose law depends on the solution
itself. \ Finally Appendix~\ref{appendix:besov} reviews some reference
material on Besov spaces and proves some technical lemmas.

\

\paragraph{Notations}We will denote $\CC^{\alpha} \assign B^{\alpha}_{\infty,
\infty} (\mathbbm{T}^2)$ the Zygmund space of regularity $\alpha \in
\mathbbm{R}$ on the torus $\mathbbm{T}^2$. See Appendix~\ref{appendix:besov}
for the definition of the general Besov spaces $B^{\alpha}_{p, q}$, the
Littlewoord--Paley operators $(\Delta_i)_{i \geqslant - 1}$ and the basic
properties thereof needed in this paper. If $V$ is a Banach space and $T > 0$,
we denote $C_T^{\alpha} V$ the space of $\alpha$--H{\"o}lder functions in $C_T
V \assign C ([0, T] ; V)$. We introduce parabolic spaces $\LL^{\alpha}_T
\assign C_T^{\alpha / 2} \CC^0 \cap C_T \CC^{\alpha}$ with norm
\begin{equation}
  \| f \|_{\LL^{\alpha}_T} = \| f \|_{C_T^{\alpha / 2} \CC^0} + \| f \|_{C_T
  \CC^{\alpha}} . \label{e:def-Lnorm}
\end{equation}
Moreover for convenience we denote $\CC_T^{\alpha} \assign C_T \CC^{\alpha}$.
We will avoid to note explicitly the time span $T$ whenever this does not
cause ambiguities. We will need also spaces for functions of $(\eta, t, x)$
where $\eta$ is an additional parameter in $[\lambda, 1]$ for $\lambda < 1$
which we denote $C_{\eta}^k V$ with norm
\begin{equation}
  \| F \|_{C_{\eta}^k V} = \sup_{\eta \in [\lambda, 1]} \sup_{n = 0, \ldots,
  k} \| \partial_{\eta}^n F (\eta,) \|_V, \label{e:def-Cetanorm}
\end{equation}
where $V$ is a Banach space of function on $[0, T] \times \mathbbm{T}^2$, in
our case $V = \CC_T^{\alpha}$ or $V = \LL^{\alpha}_T$.

We will denote by $K_{i, x} (y) = 2^{2 i} K (2^i (x - y))$ the kernel of the
Littlewood--Paley operator $\Delta_i$ and $Q_{i, t} (s) = Q_i (t - s) = 2^{2
i} Q (2^{2 i} (t - s))$ is a smoothing kernel at scale $2^{2 i}$ in the time
direction where $Q$ is a smooth, positive function with compact support in
$\mathbbm{R}_+$ and mass $1$. We introduce also the shortcut $P_i = K_{< i -
1}$. Another notation shortcut widely used in this article is to write
$\int_{x, y}$ for integrals on $\mathbbm{T}^2$ or $\mathbbm{R}$ with respect
to the measures $\mathd x$ and $\mathd y$ without specifying the integration
bounds, whenever this does not create ambiguity. Finally, we will note $\delta
f^{t x}_{s y} = f (t, x) - f (s, y)$ and $\delta_{\tau} f_{s y}^{t x} = f (s,
y) + \tau (f (t, x) - f (s, y))$ for $\tau \in [0, 1]$.

\section{Nonlinear paraproducts}\label{s:nonlinear}

In this section we introduce the nonlinear paraproduct and related results
that will be used in Section~\ref{s:paracontrolled} to analyse
equation~(\ref{eq:ur-formulation}).

Let $g : [0, T] \times \mathbbm{T}^2 \rightarrow \mathbbm{R}$, and $h :
\mathbbm{R} \times [0, T] \times \mathbbm{T}^2 \rightarrow \mathbbm{R}$ be
smooth functions. We can decompose the composition $h (g (\cdummy), \cdummy)$
via \tmtextit{nonlinear} paraproducts as follows. Define
\begin{eqnarray}
  \Pi_{\prec} (g, h) (t, x) & \assign & \sum_q \int_{y, z} P_{q, x} (y) K_{q,
  x} (z) h (g (t, y), t, z)  \label{e:nonlin-par-small}\\
  \Pi_{\circ} (g, h) (t, x) & \assign & \sum_{k \sim q} \int_{y, z} K_{k, x}
  (y) K_{q, x} (z) h (g (t, y), t, z)  \label{e:nonlin-par-same}\\
  \Pi_{\succ} (g, h) (t, x) & \assign & \sum_k \int_{y, z} K_{k, x} (z) P_{k,
  x} (y) h (g (t, z), t, y) .  \label{e:nonlin-par-big}
\end{eqnarray}
This gives a map
\begin{equation}
  (g, h) \mapsto \Pi_{\nocomma \diamondsuit} (g, h) \assign \Pi_{\prec} (g, h)
  + \Pi_{\circ} (g, h) + \Pi_{\succ} (g, h) = h (g (\cdot), \cdot) 
  \label{e:nonlin-decomp}
\end{equation}
that can be uniquely extended to
\[ \Pi_{\diamondsuit} : \CC_T^{\rho} \times C^2_{\eta} \CC_T^{\gamma}
   \rightarrow \CC_T^{\gamma \wedge \rho} \qquad \rho \in (0, 1), \gamma \in
   \mathbbm{R}, \rho + \gamma > 0, \]
thanks to the following bounds:

\begin{lemma}[Nonlinear paraproduct estimates]
  \label{l:decomp-est} Let $g \in \CC_T^{\rho}$ for some $\rho \in (0, 1)$,
  and $h \in C^2_{\eta} \CC_T^{\gamma}$ for any $\gamma \in \mathbbm{R}$. Then
  \[ \| \Pi_{\prec} (g, h) \|_{\CC^{\gamma}_T} \lesssim \| h \|_{C_{\eta}
     \CC^{\gamma}_T}, \qquad \| \Pi_{\succ} (g, h) \|_{\CC^{\rho \wedge (\rho
     + \gamma)}_T} \lesssim \| g \|_{\CC^{\rho}_T} \| h \|_{C_{\eta}^1
     \CC^{\gamma}_T}, \]
  and
  \[ \| \Pi_{\prec} (g_1, h) - \Pi_{\prec} (g_2, h) \|_{\CC^{\gamma}_T}
     \lesssim \| g_1 - g_2 \|_{C_T L^{\infty}} \| h \|_{C_{\eta}^1
     \CC^{\gamma}_T}, \qquad \]
  \begin{eqnarray*}
    \| \Pi_{\succ} (g_1, h) - \Pi_{\succ} (g_2, h) \|_{\CC_T^{\rho \wedge
    (\rho + \gamma)}} & \lesssim & \| g_1 - g_2 \|_{C_T L^{\infty}} \left( \|
    g_1 \|_{\CC^{\rho}_T} + \| g_2 \|_{\CC^{\rho}_T} \right) \| h
    \|_{C_{\eta}^2 \CC^{\gamma}_T}\\
    &  & + \| g_1 - g_2 \|_{\CC^{\rho}_T} \| h \|_{C_{\eta}^1 \CC^{\gamma}_T}
    .
  \end{eqnarray*}
  Moreover if $\rho + \gamma > 0$ we have also
  \[ \| \Pi_{\circ} (g, h) \|_{\CC^{\rho + \gamma}_T} \lesssim \| h
     \|_{C_{\eta}^1 \CC^{\gamma}_T} \| g \|_{\CC^{\rho}_T}, \]
  \begin{eqnarray*}
    \| \Pi_{\circ} (g_1, h) - \Pi_{\circ} (g_2, h) \|_{\CC_T^{\rho + \gamma}}
    & \lesssim & \| g_1 - g_2 \|_{C_T L^{\infty}} \left( \| g_1
    \|_{\CC^{\rho}_T} + \| g_2 \|_{\CC^{\rho}_T} \right) \| h \|_{C_{\eta}^2
    \CC^{\gamma}_T}\\
    &  & + \| g_1 - g_2 \|_{\CC^{\rho}_T} \| h \|_{C_{\eta}^1 \CC^{\gamma}_T}
    .
  \end{eqnarray*}
  In particular if $\rho + \gamma > 0$ the composition $\Pi_{\diamondsuit} (g,
  h) = h (g (\cdot), \cdot)$ is linear in $h$ and locally Lipshitz in $g$:
  \[ \| \Pi_{\diamondsuit} (g, h) \|_{\CC_T^{\gamma}} \lesssim \| h
     \|_{C_{\eta}^1 \CC^{\gamma}_T} \| g \|_{\CC^{\rho}_T}, \]
  \[ \| \Pi_{\diamondsuit} (g_1, h) - \Pi_{\diamondsuit} (g_2, h)
     \|_{\CC_T^{\gamma}} \lesssim \| g_1 - g_2 \|_{\CC^{\rho}_T} \left( 1 + \|
     g_1 \|_{\CC^{\rho}_T} + \| g_2 \|_{\CC^{\rho}_T} \right) \| h
     \|_{C_{\eta}^2 \CC^{\gamma}} . \]
\end{lemma}

\begin{proof}
  Using the fact that
  \[ \| \Delta_k h (g (t, y), t, \cdot) \|_{L^{\infty}} \lesssim 2^{- \gamma
     k} \| h \|_{C_{\eta} \CC_T^{\gamma}}, \]
  \[ \| \Delta_k h (g (t, y), t, \cdot) - \Delta_k h (g (t, y'), t, \cdot)
     \|_{L^{\infty}} \lesssim 2^{- \gamma k} \| h \|_{C_{\eta}^1
     \CC_T^{\gamma}} \| g \|_{C_T  \CC^{\rho}} | y - y' |^{\rho}, \]
  and
  \[ \| \Delta_k h (g_1 (t, y), t, \cdot) - \Delta_k h (g_2 (t, y), t, \cdot)
     \|_{L^{\infty}} \lesssim 2^{- \gamma k} \| h \|_{C_{\eta}^1
     \CC_T^{\gamma}} \| g_1 - g_2 \|_{C_T L^{\infty}}, \]
  we obtain the bounds on $\Pi_{\prec} (g, h)$, $\Pi_{\succ} (g, h)$,
  $\Pi_{\circ} (g, h)$ and $\Pi_{\prec} (g_1, h) - \Pi_{\prec} (g_2, h)$. We
  proceed to estimate the term $\Pi_{\succ} (g_1, h) - \Pi_{\succ} (g_2, h)$.
  We will use the following notation for brevity:
  \[ \delta g_{2 z}^{1 y} : = g_1 (t, y) - g_2 (t, z) \quad \tmop{and} \quad
     \delta_{\tau} g_{2 z}^{1 y} : = g_2 (t, z) + \tau (g_1  (t, y) - g_2 (t,
     z)) . \]
  Then
  \begin{eqnarray*}
    &  & \left| \int_{y, z} K_{k, x} (z) P_{k, x} (y) [h (g_1 (t, z), t, y) -
    h (g_2 (t, z), t, y)] \right|\\
    & = & \left| \int_{y, z, \tau \in [0, 1]} K_{k, x} (z) P_{k, x} (y)
    [\partial_{\eta} h (\delta_{\tau} g_{2 z}^{1 z}, t, y) \delta g_{2 z}^{1
    z} - \partial_{\eta} h (\delta_{\tau} g_{2 x}^{1 x}, t, y) \delta g_{2
    x}^{1 x}] \right|\\
    & \lesssim & \left| \int_{y, z, t \in [0, 1], \sigma \in [0, 1]} K_{k, x}
    (z) P_{k, x} (y) \partial^2_{\eta} h (\delta_{\sigma} (\delta_{\tau}
    g_2^1)^z_x, t, y) (\delta g^{1 z}_{1 x} - \delta g^{2 z}_{2 x}) \delta
    g_{2 z}^{1 z} \right|\\
    &  & + \left| \int_{y, z, \tau \in [0, 1]} K_{k, x} (z) P_{k, x} (y)
    \partial_{\eta} h (\delta_{\tau} g_{2 x}^{1 x}, t, y) (\delta g_{2 z}^{1
    z} - \delta g_{2 x}^{1 x}) \right|\\
    & \lesssim & \| g_1 - g_2 \|_{C_T L^{\infty}} \left( \| g_1 \|_{C_T
    \CC^{\rho}} + \| g_2 \|_{C_T \CC^{\rho}} \right) \| h \|_{C_{\eta}^2 C_T
    \CC^{\gamma}} 2^{- \rho k} \sum_{q < k - 1} 2^{- \gamma q}\\
    &  & + \| g_1 - g_2 \|_{C_T \CC^{\rho}}  \| h \|_{C_{\eta}^1 C_T
    \CC^{\gamma}} 2^{- \rho k} \sum_{q < k - 1} 2^{- \gamma q} .
  \end{eqnarray*}
  With the same reasoning we can bound the norm of $\Pi_{\circ} (g_1, h) -
  \Pi_{\circ} (g_2, h)$.
\end{proof}

We will need a time-smoothed nonlinear paraproduct, defined as
\begin{equation}
  \Pi_{\precprec} (g, h) (t, x) \assign \sum_i \int_{y, s} Q_{i, t} (s) P_{i,
  x} (y) (\Delta_i h (g (s, y), t, \cdot)) (x), \label{e:def-nonlinear}
\end{equation}
where, by convention, a continuous function $t \mapsto g (t)$ on
$\mathbbm{R}_+$ is extended to $\mathbbm{R}$ by defining $g (t) = g (0)$ for
$t \leqslant 0$. This convention preserves the H{\"o}lder norms of index in
$[0, 1]$. The modified nonlinear paraproduct enjoys similar bounds to the
regular one.

\begin{lemma}
  \label{l:nonlinear-est}Let $g \in C_T L^{\infty}$ and $h \in C^1_{\eta}
  \LL_T^{\gamma}$ for $\gamma \in (0, 2)$. Then
  \[ \| \Pi_{\precprec} (g, h) \|_{\CC^{\gamma}_T} \lesssim \| h \|_{C_{\eta}
     \CC_T^{\gamma}} \quad \tmop{and} \quad \| \Pi_{\precprec} (g, h)
     \|_{\LL^{\gamma}_T} \lesssim \| h \|_{C_{\eta} \LL_T^{\gamma}} . \]
  Moreover, $\Pi_{\precprec} (g, h)$ is linear in $h$ and:
  \[ \| \Pi_{\precprec} (g_1, h) - \Pi_{\precprec} (g_2, h)
     \|_{\CC_T^{\gamma}} \lesssim \| g_1 - g_2 \|_{C_T L^{\infty}} \| h
     \|_{C^1_{\eta} \CC_T^{\gamma}}, \]
  \[ \| \Pi_{\precprec} (g_1, h) - \Pi_{\precprec} (g_2, h)
     \|_{\LL_T^{\gamma}} \lesssim \| g_1 - g_2 \|_{C_T L^{\infty}} \| h
     \|_{C^1_{\eta} \LL_T^{\gamma}} . \]
\end{lemma}

\begin{proof}
  The norm~$\| \Pi_{\precprec} (g, h) \|_{C_T \CC^{\gamma}}$ can be treated in
  the same way as in Lemma~\ref{l:decomp-est}. We estimate~$\| \Pi_{\precprec}
  (g, h) \|_{C_T^{\gamma / 2} \CC^0}$ as follows:
  \begin{eqnarray*}
    &  & \| \Delta_j \Pi_{\precprec} (g, h) (t_1) - \Delta_j \Pi_{\precprec}
    (g, h) (t_2) \|_{L^{\infty}}\\
    & \lesssim & \sup_x \left| \int_z K_{j, x} (z) \sum_{i \sim j} \int_{y,
    s} Q_{i, t_1} (s) P_{i, z} (y) [\Delta_i h (g (s, y), t_1, z) - \Delta_i h
    (g (s, y), t_2, z)] \right|\\
    &  & + \sup_x \left| \int_z K_{j, x} (z) \sum_{i \sim j} \int_{y, s}
    [Q_{i, t_1} (s) - Q_{i, t_2} (s)] P_{i, z} (y) \Delta_i h (g (s, y), t_2,
    z) \right|\\
    & \lesssim & \| h (\cdummy, t_1, \cdummy) - h (\cdummy, t_2, \cdummy)
    \|_{C_{\eta} \CC^0} + | t_1 - t_2 |^{\gamma / 2} \| h \|_{C_{\eta}
    \CC_T^{\gamma}} .
  \end{eqnarray*}
  The second inequality can be obtained easily with the same techniques used
  so far.
\end{proof}

\subsection{Nonlinear commutator}

The next technical ingredient is a commutator lemma between the non-linear
paraproduct of (\ref{e:def-nonlinear}) and the standard resonant product. It
will be needed below to analyse a term of the form $\Pi_{\precprec} (g, h)
\circ \Delta \Pi_{\precprec} (g, h)$ so we will specialise our discussion to
this specific structure. Notice that in the following the various space--time
operators act pointwise in the parameter $\eta$, in the sense that, for
example:
\[ (h \circ \Delta h) (\eta, t, x) = (h (\eta, t, \cdot) \circ \Delta h (\eta,
   t, \cdot)) (x) . \]
\begin{lemma}
  \label{l:nonlin-comm}We define the map $\Lambda : C^{\infty} ([0, T],
  \mathbbm{T}^2) \times C_{\eta}^2 C^{\infty} ([0, T], \mathbbm{T}^2)
  \rightarrow C^{\infty} ([0, T], \mathbbm{T}^2)$ by
  \[ \Lambda (g, h) \assign [\Pi_{\precprec} (g, h) \circ \Delta
     \Pi_{\precprec} (g, h)] - \Pi_{\diamondsuit} (g, h \circ \Delta h) . \]
  Then for all \ $\rho \in (0, 1)$, $\gamma <\mathbbm{R}$, $\varepsilon > 0$
  such that $2 \gamma - 2 + \rho - \varepsilon > 0$, we have
  \[ \| \Lambda (g, h) \|_{\CC_T^{2 \gamma - 2 + \rho - \varepsilon}}
     \lesssim \| g \|_{\LL_T^{\rho}} \| h \|_{C^1_{\eta} \CC^{\gamma}_T}^2 \]
  and
  \begin{eqnarray*}
    \| \Lambda (g_1, h) - \Lambda (g_2, h) \|_{\CC_T^{2 \gamma - 2 + \rho -
    \varepsilon}} & \lesssim & \| g_1 - g_2 \|_{C_T L^{\infty}} \left( \| g_1
    \|_{\LL^{\rho}_T} + \| g_2 \|_{\LL^{\rho}_T} \right) \| h \|^2_{C_{\eta}^2
    \CC^{\gamma}_T}\\
    &  & + \| g_1 - g_2 \|_{\LL^{\rho}_T} \| h \|^2_{C_{\eta}^1
    \CC^{\gamma}_T} .
  \end{eqnarray*}
  As a consequence $\Lambda$ can be uniquely extended to a locally Lipshitz
  function
  \[ \Lambda : \LL^{\rho}_T \times C_{\eta}^2 \CC_T^{\gamma} \rightarrow
     \CC_T^{2 \gamma - 2 + \rho - \varepsilon} . \]
\end{lemma}

\begin{proof}
  We let $\eta (t, x) \assign \Pi_{\precprec} (g, h) \circ \Delta
  \Pi_{\precprec} (g, h) (t, x)$ . Then
  \begin{eqnarray}
    \Delta_q \eta (t, x) & = & \int_z K_{q, x} (z) (\Pi_{\precprec} (g, h)
    \circ \Delta \Pi_{\precprec} (g, h)) (t, z) \nonumber\\
    & = & \int_z K_{q, x} (z) (\Pi_{\precprec} (g (t, z), h) \circ \Delta
    \Pi_{\precprec} (g (t, z), h)) (t, z) \nonumber\\
    &  & + \int_z K_{q, x} (z) ((\Pi_{\precprec} (g, h) - \Pi_{\precprec} (g
    (t, z), h)) \circ \Delta \Pi_{\precprec} (g (t, z), h)) (t, z) 
    \label{e:Deltaeta2}\\
    &  & + \int_z K_{q, x} (z) (\Pi_{\precprec} (g, h) \circ \Delta
    (\Pi_{\precprec} (g, h) - \Pi_{\precprec} (g (t, z), h))) (t, z) 
    \label{e:Deltaeta3}
  \end{eqnarray}
  Now notice that
  \[ \Delta_j [\Delta \Pi_{\precprec} (g (t, z), h)] (t, z) = \Delta_j
     \Pi_{\precprec} (g (t, z), \Delta h) (t, z) \]
  and then
  \begin{eqnarray*}
    (\Pi_{\precprec} (g (t, z), h) \circ \Pi_{\precprec} (g (t, z), \Delta h))
    (t, z) & = & \sum_{i \sim j} \Delta_i h (g (t, z), t, \cdummy) (z)
    \Delta_j [\Delta h (g (t, z), t, \cdummy)] (z)\\
    & = & \Pi_{\diamondsuit} (g, h \circ \Delta h) (t, z)
  \end{eqnarray*}
  We proceed therefore to estimate (\ref{e:Deltaeta2}) and
  (\ref{e:Deltaeta3}). We obtain
  \[ \int_z K_{q, x} (z) [(\Pi_{\precprec} (g, h) - \Pi_{\precprec} (g (t, z),
     h)) \circ \Delta \Pi_{\precprec} (g (t, z), h)] (t, z) \]
  \[ = \int_z K_{q, x} (z) \sum_{i \sim j \gtrsim q} (\Delta_i \Pi_{\precprec}
     (g, h) (t, z) - \Delta_i \Pi_{\precprec} (g (t, z), h) (t, z)) \Delta_j
     \Delta \Pi_{\precprec} (g (t, z), h) (t, z) . \]
  Using Lemma~\ref{l:nonlinear-est} we have
  \[ | \Delta \Delta_j \Pi_{\precprec} (g (t, z), h) (t, z) | \lesssim 2^{(2 -
     \gamma) j} \| h \|_{C_{\eta} \CC^{\gamma}_T} . \]
  Lemma~\ref{l:para-increment} gives
  \[ | \Delta_i (\Pi_{\precprec} (g, h) - \Pi_{\precprec} (g (t, z), h)) (t,
     z) | \lesssim 2^{- (\gamma + \rho - \varepsilon) i} \| g
     \|_{\LL^{\rho}_T} \| h \|_{C_{\eta}^1 \CC^{\gamma}_{_T}}, \]
  and thus (\ref{e:Deltaeta2}) is bounded by $2^{- (\gamma + \rho - 2 -
  \varepsilon) q} \| g \|_{\LL^{\rho}_T} \| h \|^2_{\CC^{\gamma}_{_T}}$.
  
  We can easily bound (\ref{e:Deltaeta3}) in the same way, and this proves the
  first inequality. For the second inequality, Lemma~\ref{l:nonlinear-est}
  yields
  \[ | \Delta_j \Delta \Pi_{\precprec} (g_1, h) (t, z) - \Delta_j \Delta
     \Pi_{\precprec} (g_2, h) (t, z) | \lesssim 2^{(2 - \gamma) j} \| g_1 -
     g_2 \|_{C_T L^{\infty}} \| h \|_{C^1_{\eta} \CC^{\gamma}_T}, \]
  and using the second inequality of Lemma~\ref{l:para-increment} we obtain
  the desired bound.
  
  The extension of $\Lambda$ to $\LL^{\rho}_T \times C_{\eta}^2
  \LL_T^{\gamma}$ is standard (see e.g. the proof of the commutator lemma
  in~{\cite{gubinelli_paracontrolled_2015}}, Lemma~2.4). 
\end{proof}

\begin{lemma}
  \label{l:para-increment}With the same assumptions of
  Lemma~\ref{l:nonlin-comm} we have
  \[ | \Delta_i \Pi_{\precprec} (g, h) (t, z) - \Delta_i \Pi_{\precprec} (g
     (t, z), h) (t, z) | \lesssim 2^{(\varepsilon - \rho - \gamma) i} \| g
     \|_{\LL^{\rho}_T} \| h \|_{C_{\eta}^1 \CC^{\gamma}_{_T}} \]
  and
  \[ \left| \Delta_i \Pi_{\precprec} (g_1, h) (t, z) - \Delta_i
     \Pi_{\precprec} (g_1 (t, z), h) (t, z) - \Delta_i \Pi_{\precprec} (g_2,
     h) (t, z) + \Delta_i \Pi_{\precprec} (g_2 (t, z), h) (t, z) \right| \]
  \[ \lesssim 2^{(\varepsilon - \rho - \gamma) i} \left[ \| g_1 - g_2
     \|_{\LL^{\rho}_T} \| h \|_{C_{\eta}^1 \CC^{\gamma}_T} + \| g_1 - g_2
     \|_{C_T L^{\infty}} \left( \| g_1 \|_{\LL^{\rho}_T} + \| g_2
     \|_{\LL^{\rho}_T} \right) \| h \|_{C_{\eta}^2 \CC^{\gamma}_T} \right] \]
\end{lemma}

\begin{proof}
  
  \begin{eqnarray*}
    &  & [\Delta_i \Pi_{\precprec} (g, h) - \Delta_i \Pi_{\precprec} (g (t,
    z), h)] (t, z)\\
    & = & \sum_{\underset{}{k \sim i}} \int_{_{x, y, s, \tau}} K_{i, z} (x)
    Q_{k, t} (s) P_{k - 1, x} (y) \partial_{\eta} \Delta_k h (\delta_{\tau}
    g^{s y}_{t z}, t, x) (\delta g^{s y}_{t y} + \delta g^{t y}_{t z})\\
    & \lesssim & \sum_{\underset{}{k \sim i}} \int_{_{x, y, s, \tau}} | K_{i,
    z} (x) Q_{k, t} (s) P_{k - 1, x} (y) | \| \partial_{\eta} \Delta_k h
    \|_{C_T L^{\infty}} | t - s |^{(\rho - \varepsilon) / 2} \| g \|_{C^{\rho
    / 2 - \varepsilon / 2}_T L^{\infty}}\\
    &  & + \sum_{\underset{}{k \sim i}} \int_{_{x, y, s, \tau}} | K_{i, z}
    (x) Q_{k, t} (s) P_{k - 1, x} (y) | \| \partial_{\eta} \Delta_k h \|_{C_T
    L^{\infty}} | y - z | \| g \|_{\CC^{\rho}_T}\\
    & \lesssim & 2^{- (\rho - \varepsilon) i} 2^{- \gamma i} \|
    \partial_{\eta} h \|_{\CC^{\gamma}_T} \left( \| g \|_{C^{\rho}_T \CC^0} +
    \| g \|_{\CC^{\rho}_T} \right)
  \end{eqnarray*}
  where we used the notation $\delta g^{s y}_{t z} = g (s, y) - g (t, z)$,
  $\delta_{\tau} g^{s y}_{t z} = g (t, z) + \tau (g (s, y) - g (t, z))$ and
  Lemma~\ref{l:interpol}. This proves the first bound.
  
  The second inequality can be obtained in the same way with the techniques
  already used here and in Lemma~\ref{l:decomp-est}.
\end{proof}

\subsection{Approximate paradifferential problem}\label{s:approx-paradiff}

In this section we construct an approximate solution to the equation
\begin{equation}
  (\partial_t - g \prec \Delta) u = f, \qquad u (0, \cdot) = 0,
  \label{e:paradiff}
\end{equation}
with data $f \in \CC^{\gamma - 2}$ and $g \in \LL^{\rho}_T$, for some fixed
$\rho, \gamma \in (0, 1)$. The idea is to obtain it via a certain class of
paradifferential operators. We introduce the operator $\LL$ acting on
functions of~$(\eta, t, x)$ by
\begin{equation}
  \left( \LL U \right) (\eta, t, x) \assign \partial_t U (\eta, t, x) - \eta
  \Delta U (\eta, t, x) \label{e:def-Loper}
\end{equation}
Observe that if $u$ does not depend on $\eta$ we can define
\begin{equation}
  \Pi_{\prec} \left( g, \LL \right) u \assign \Pi_{\prec} \left( g, \LL u
  \right) . \label{e:paradiff-oper}
\end{equation}
and we obtain $\Pi_{\prec} \left( g, \LL \right) u = \partial_t u - g \prec
\Delta u$.

We can describe the commutation between the differential $\LL$ and the
paraproduct $\Pi_{\precprec} (g, \cdummy)$ via the following estimate:

\begin{lemma}
  \label{l:def-psi}Let $\rho \in (0, 1)$, $\gamma \in \mathbbm{R}$. Let $U \in
  C_{\eta}^2 \CC^{\gamma}_T$ and $g \in \LL^{\rho}_T$ such that $g \in
  [\lambda, 1]$. Define
  \[ \Psi (g, U) \assign R_1 + R_2 \]
  with $R_1$ and $R_2$ as in (\ref{e:rem-1}),~(\ref{e:rem-2}). Then for every
  $\varepsilon > 0$
  
  
  \begin{equation}
    \| \Psi (g, U) \|_{\CC_T^{\rho + \gamma - 2 - \varepsilon}} \lesssim (1 +
    \| g \|_{C_T L^{\infty}}) \| g \|_{\LL^{\rho}_T}^{} \| U \|_{C_{\eta}^1
    \CC^{\gamma}_{_T}}  \label{e:psi-est} .
  \end{equation}
  Moreover, $\Psi (g, U)$ is linear in $U$ and
  \[ \| \Psi (g_1, U) - \Psi (g_2, U) \|_{\CC_T^{\rho + \gamma - 2 -
     \varepsilon}} \lesssim \| g_1 - g_2 \|_{\LL^{\rho}_T} \left( 1 + \| g_1
     \|_{\LL^{\rho}_T} + \| g_2 \|_{\LL^{\rho}_T} \right) \| U \|_{C_{\eta}^2
     \CC^{\gamma}_{_T}} . \]
  In particular, we have
  \begin{equation}
    \Psi (g, U) = \Pi_{\precprec} \left( g, \LL U \right) - \Pi_{\prec} \left(
    g, \LL \right) \Pi_{\precprec} (g, U) \in C_T \CC^{\rho + \gamma - 2 -
    \varepsilon}  \label{e:quasidef-psi}
  \end{equation}
  whenever this expression makes sense.
\end{lemma}

\begin{proof}
  We start considering $g \in C^{\infty} ([0, T], \mathbbm{T}^2)$ and $U \in
  C_{\eta}^2 C^{\infty} ([0, T], \mathbbm{T}^2)$, and
  prove~$(\ref{e:quasidef-psi})$ in this setting. Notice that $\Pi_{\precprec}
  \left( g (t, y), \LL_{g (t, y)} U \right) = \LL U (g (t, y))$. As a
  consequence, we can estimate
  \begin{eqnarray*}
    &  & \Pi_{\precprec} \left( g, \LL U \right) (t, x) - \Pi_{\prec} \left(
    g, \LL \Pi_{\precprec} (g, U) \right) (t, x)\\
    & = & \Pi_{\precprec} \left( g, \LL U \right) (t, x) - \sum_k \int_y
    P_{k, x} (y) \left( \LL_{g (t, y)} \Delta_k \Pi_{\precprec} (g, U) \right)
    (t, x)\\
    & = & \Pi_{\precprec} \left( g, \LL U \right) (t, x) - \sum_k \int_y
    P_{k, x} (y) (\partial_t \Delta_k \Pi_{\precprec} (g, U)) (t, x)  \\
    &  & + \sum_k \int_y P_{k, x} (y) g (t, y) (\Delta \Delta_k
    \Pi_{\precprec} (g, U)) (t, x)\\
    & = & \Pi_{\precprec} \left( g, \LL U - \partial_t U \right) (t, x) +
    \sum_k \int_y P_{k, x} (y) g (t, y) (\Delta_k \Pi_{\precprec} (g, \Delta
    U)) (t, x)\\
    &  & + \sum_k \int_{} P_{k, x} (y) g (t, y) (\Delta_k [\Delta,
    \Pi_{\precprec} (g, \cdot)] U) (t, x) - \sum_k \int_{} P_{k, x} (y)
    (\Delta_k [\partial_t, \Pi_{\precprec} (g, \cdot)] U) (t, x)
  \end{eqnarray*}
  We have
  \begin{eqnarray*}
    \Pi_{\precprec} \left( g, \LL U - \partial_t U \right) (t, x) + \sum_k
    \int_y P_{k, x} (y) g (t, y) (\Delta_k \Pi_{\precprec} (g, \Delta U)) (t,
    x) & = & R_1 (t, x)
  \end{eqnarray*}
  with the definition
%  \[ R_1 (t, x) \assign \]
  \begin{equation}
    R_1 (t, x) \assign  \sum_{k, i} \int_{_{\overset{y', s}{y, z}}} P_{k, x} (y) K_{k, x} (z)
    P_{i, z} (y') Q_{i, t} (s) [g (t, y) - g (s, y')] \Delta \Delta_i U (g (s,
    y'), t, z) \label{e:rem-1}
  \end{equation}
  and
  \[ \sum_k \int_y P_{k, x} (y) [g (t, y) (\Delta_k [\Delta, \Pi_{\precprec}
     (g, \cdot)] U) (t, x) - (\Delta_k [\partial_t, \Pi_{\precprec} (g,
     \cdot)] U) (t, x)] = R_2 (t, x) \]
  with the definition
  \begin{eqnarray}
    R_2 (t, x) & \assign & \sum_{k, i} \int_{y, y', s} P_{k, x} (y) K_{k, x}
    (z) Q_{i, t} (s) g (t, y) \Delta P_{i, z} (y') \Delta_i U (g (s, y'), t,
    z) \nonumber\\
    &  & + 2 \sum_{k, i} \int_{y, y', s} P_{k, x} (y) K_{k, x} (z) Q_{i, t}
    (s) g (t, y) \nabla P_{i, z} (y') \nabla \Delta_i U (g (s, y'), t, z)
    \nonumber\\
    &  & - \sum_{k, i} \int_{y, y', s} P_{k, x} (y) K_{k, x} (z) \partial_t
    Q_{i, t} (s) P_{i, z} (y') \Delta_i U (g (s, y'), t, z) .  \label{e:rem-2}
  \end{eqnarray}
  Indeed:
  \begin{eqnarray}
    ([\partial_t, \Pi_{\precprec} (g, \cdot)] U) (t, x) & = & \sum_i \int_{y,
    s} (\partial_t Q_{i, t}) (s) P_{i, x} (y) (\Delta_i U (g (s, y), t, x)),
    \nonumber\\
    ([\Delta, \Pi_{\precprec} (g, \cdot)] U) (t, x) & = & \sum_i \int_{y, s}
    Q_{i, t} (s) \Delta P_{i, x} (y) (\Delta_i U (g (s, y), t, x)) 
    \label{eq:q-time}\\
    &  & + 2 \sum_i \int_{y, s} Q_{i, t} (s) \nabla P_{i, x} (y) (\nabla
    \Delta_i U (g (s, y), t, x)) . \nonumber
  \end{eqnarray}
  This shows that $(\ref{e:quasidef-psi})$ holds for smooth functions.
  
  With the techniques used in Lemma~\ref{l:para-increment} we can estimate
  \[ | \Delta_q R_1 (t, x) | \lesssim \sum_{k \sim q} \left( 2^{- (\rho -
     \varepsilon) k} \| g \|_{C_T^{\rho / 2} \CC^0} + 2^{- \rho k} \| g
     \|_{\CC^{\rho}_T} \right) 2^{(2 - \gamma) k} \| U \|_{C_{\eta}
     \CC^{\gamma}_T} . \]
  By the spectral support properties of the commutators we have that
  \begin{eqnarray*}
    \| [\Delta, \Pi_{\precprec} (g, \cdot)] U \|_{\CC_T^{\gamma + \rho - 2}} &
    \lesssim & \| g \|_{\CC^{\rho}_T} \| U \|_{C^1_{\eta} \CC^{\gamma}_T},\\
    \| \Delta_q [\partial_t, \Pi_{\precprec} (g, \cdot)] U \|_{C_T L^{\infty}}
    & \lesssim & \left( 2^{(2 + \varepsilon - \rho - \gamma) q} \| g
    \|_{C_T^{\rho / 2} \CC^0} + 2^{(2 - \rho - \gamma) q} \| g
    \|_{\CC^{\rho}_T} \right) \| U \|_{C^1_{\eta} \CC^{\gamma}_T} .
  \end{eqnarray*}
  This yields
  \[ \| R_2 \|_{\CC_T^{\gamma + \rho - 2 - \varepsilon}} \lesssim (1 + \| g
     \|_{C_T L^{\infty}}) \| g \|_{\LL^{\rho}_T} \| U \|_{C^1_{\eta}
     \CC^{\gamma}_T} . \]
  We have so far proved (\ref{e:psi-est}) and then (\ref{e:quasidef-psi})
  follows by continuity. The local Lipshitz dependence on $g$ can be obtained
  via similar computations.
\end{proof}

\begin{remark}
  If $f$ does not depends on $\eta$ we consider the parametric problem
  \begin{equation}
    (\partial_t - \eta \Delta) U_f (\eta, t) = f, \qquad U_f (\eta, 0) = 0,
    \qquad \eta \in [\lambda, 1], \label{e:def-U}
  \end{equation}
  which is solved by
  \[ U_f (\eta, t) = \int_0^t e^{\eta \Delta (t - s)} f \mathd s. \]
  Remark that
  \[ \partial_{\eta} U_f (\eta, t) = \int_0^t e^{\eta \Delta (t - s)} (t - s)
     \Delta f \mathd s \quad \tmop{and} \quad \partial_{\eta}^2 U_f (\eta, t)
     = \int_0^t e^{\eta \Delta (t - s)} (t - s)^2 \Delta^2 f \mathd s \]
  so that we have (since $\eta \geqslant \lambda$):
  \begin{equation}
    \| U_f \|_{C_{\eta}^2 \LL^{\gamma}_T} \assign \sup_{n = 0, 1, 2}
    \sup_{\eta \in [\lambda, 1]} \| \partial_{\eta}^n U_f (\eta)
    \|_{\LL_T^{\gamma}} \lesssim \| f \|_{\CC^{\gamma - 2}_T}  \label{e:U-est}
  \end{equation}
  We define then
  \begin{equation}
    u (t, x) \assign \Pi_{\precprec} (g, U_f) (t, x) \label{e:def-u}
  \end{equation}
  and observe that $u (t, x)$ is an approximate solution of
  equation~(\ref{e:paradiff}), indeed
  \[ (\partial_t - g \prec \Delta) u = \Pi_{\prec} \left( g, \LL
     \Pi_{\precprec} (g, U_f) \right) = \Pi_{\precprec} \left( g, \LL U_f
     \right) - \Psi (g, U_f) = f - \Psi (g, U_f) \]
  and the estimation in Lemma~\ref{l:def-psi} together with the
  bound~(\ref{e:U-est}) yield immediately the following inequality:
  \begin{equation}
    \| \Psi (g, U_f) \|_{\CC^{\rho + \gamma - 2 - \varepsilon}_T} \lesssim \|
    g \|_{\LL^{\rho}_T} (1 + \| g \|_{C_T L^{\infty}}) \| f \|_{\CC^{\gamma -
    2}_T} . \label{e:psi-est2}
  \end{equation}
\end{remark}

\section{Paracontrolled Ansatz\label{s:paracontrolled}}

In order to give a meaning to the PDE in (\ref{e:paracontrolled-eq}) with
initial condition $u_0 \in \CC^{\alpha}$, our initial goal will be to get
informations on solutions $\theta = \theta (g)$ of the equation
\[ \partial_t \theta - g \prec \Delta \theta = \xi, \]
for a fixed $g \in \CC^{\alpha}_T$, $2 / 3 < \alpha < 1$. Using the results of
Section~\ref{s:approx-paradiff}, we consider to this effect the parametric
problem
\[ (\partial_t - \eta \Delta) \vartheta (\eta, t) = \xi, \]
for $\eta \in [\lambda, 1]$. We will consider the stationary solution of this
problem which has the form
\begin{equation}
  \vartheta (\eta, x) = \int_0^{\infty} e^{\eta \Delta s} \xi \mathd s = (-
  \eta \Delta)^{- 1} \xi \label{eq:def-theta}
\end{equation}
and in order to have this well defined we impose that the noise $\xi$ has zero
mean on $\mathbbm{T}^2$ (this is a simplifying assumption which can be easily
removed, e.g. at the price of adding a linear term to the equation). It is
easy to see that
\begin{equation}
  \| \vartheta \|_{C_{\eta}^2 \LL^{\alpha}_T} = \| \vartheta \|_{C_{\eta}^2
  \CC^{\alpha}_T} \lesssim \| \xi \|_{\CC^{\alpha - 2}} .
  \label{eq:vartheta-est} 
\end{equation}
We define now for every $t \in [0, T]$
\[ \theta (t, x) \assign \Pi_{\precprec} (a (u), \vartheta) . \]
Thanks to Lemma~\ref{l:nonlinear-est} we have the bound $\| \theta
\|_{\LL^{\alpha}_{_T}} \lesssim \| \vartheta \|_{C_{\eta} \LL^{\alpha}_{_T}}
\lesssim \| \xi \|_{\CC^{\alpha - 2}}$ . We observe that this definition
together with Lemma~\ref{l:def-psi} gives
\[ \partial_t \theta - a (u) \prec \Delta \theta = \xi - \Psi (a (u),
   \vartheta) \]
with $\| \Psi (a (u), \vartheta) \|_{\CC^{2 \alpha - 2 - \varepsilon}_T}
\lesssim \| a (u) \|^2_{\LL^{\alpha}_T} \| \xi \|_{\CC^{\alpha - 2}_T}$ . We
expect then $\Psi (a (u), \vartheta)$ to be bounded in $\CC^{2 \alpha - 2 -
\varepsilon}_T$ for any $\varepsilon > 0$. At this point let us introduce the
Ansatz
\begin{equation}
  u = \theta + u^{\sharp} . \label{e:firstansatz}
\end{equation}
\begin{remark}
  Notice that we are not making any assumption on the existence of such~$u$,
  which is the subject of Section~\ref{s:fixpoint}. Our aim here is to find
  the equation that a couple $(u, u^{\sharp}) \in \CC_T^{\alpha} \times
  \CC_T^{2 \alpha}$ verifying (\ref{e:firstansatz}) must solve, in order for
  $u$ to solve (\ref{e:paracontrolled-eq}).
\end{remark}

Observe that
\[ \partial_t u - a (u) \prec \Delta u = (\partial_t - a (u) \prec \Delta)
   \theta + (\partial_t - a (u) \prec \Delta) u^{\sharp} = \xi + (\partial_t -
   a (u) \prec \Delta) u^{\sharp} - \Psi (a (u), \vartheta) . \]
It follows that $u^{\sharp}$ must solve
\begin{equation}
  \left\{ \begin{array}{l}
    (\partial_t - a (u) \prec \Delta) u^{\sharp} = \Phi (u) + \Psi (a (u),
    \vartheta)\\
    u^{\sharp} (t = 0) = u^{\sharp}_0 \assign u_0 - \Pi_{\precprec} (a (u_0),
    \vartheta) (t = 0) \in \CC^{\alpha}
  \end{array} \right. \label{e:usharp-equadiff} \qquad
\end{equation}
with $\Phi (u) = a (u) \circ \Delta u + a (u) \succ \Delta u$, and if we can
make sense of the resonant term $a (u) \circ \Delta u$, it is reasonable to
expect $u^{\sharp} (t, \cdummy) \in \CC^{2 \alpha}$ $\forall t \in (0, T]$.
Indeed, take $U^{\sharp} \assign U_Q$ to be the solution of
\begin{equation}
  (\partial_t - \eta \Delta) U^{\sharp} (\eta, t) = Q \qquad U^{\sharp} (\eta,
  0) = 0 \label{e:def-Usharp} 
\end{equation}
for some $Q = Q (u^{\sharp})$ to be determined and $\eta \in [\lambda, 1]$.
Using again Lemma~\ref{l:def-psi} we have
\[ (\partial_t - a (u) \prec \Delta) \Pi_{\precprec} (a (u), U^{\sharp}) = Q
   (u^{\sharp}) - \Psi (a (u), U^{\sharp}) . \]
For $\eta \in [\lambda, 1] $we define $\PP_t u_0^{\sharp} (\eta) \assign
e^{\eta \Delta t} u_0^{\sharp}$ so that $\LL \left( \PP_t u_0^{\sharp} \right)
= 0$, with $\LL$ as in (\ref{e:def-Loper}).

We set
\begin{equation}
  u^{\sharp} \assign \Pi_{\precprec} (a (u), U^{\sharp}) + \Pi_{\precprec}
  \left( a (u), \PP u_0^{\sharp} \right) . \label{eq:decomp-sharp}
\end{equation}
Taking
\[ Q (u^{\sharp}) : = \Phi (u) + \Psi (a (u), \vartheta) + \Psi (a (u),
   U^{\sharp}) + \Psi \left( a (u), \PP u_0^{\sharp} \right), \]
we obtain that $U^{\sharp}$ solves equation~(\ref{e:def-Usharp}) if and only
if $u^{\sharp}$ solves equation~(\ref{e:usharp-equadiff}). As we will see, $Q
(u^{\sharp}) (t)$ belongs to~$\CC^{2 \alpha - 2}$ $\forall t \in (0, T]$ but
not uniformly as $t \rightarrow 0$. However it belongs to $\CC^{\alpha - 2}$
uniformly as $t \rightarrow 0$.

It remains to control the resonant term $a (u) \circ \Delta u$ appearing in
$\Phi (u)$. We have
\[ a (u) \circ \Delta u = a (u) \circ \Delta \theta + a (u) \circ \Delta
   u^{\sharp} . \]
By paralinearization (see Theorem~\ref{th:paralinearisation}) $a (u) = a' (u)
\prec u + R_a (u)$ with $\| R_a (u) \|_{\CC^{2 \alpha}_T} \lesssim 1 + \| u
\|^2_{\CC^{\alpha}_T}$, and then
\[ a (u) \circ \Delta \theta = (a' (u) \prec u) \circ \Delta \theta + R_a (u)
   \circ \Delta \theta . \]
In order to use the commutator lemma (Lemma~\ref{lem:commutator}) we can
estimate $a' (u)$, recalling that $\alpha \in (0, 1)$, as
\[ \| a' (u) \|_{\CC^{\alpha}_T} \lesssim \| a'' \|_{L^{\infty}} \| u
   \|_{\CC^{\alpha}_T} \]
and write
\[ a (u) \circ \Delta \theta = a' (u) (u \circ \Delta \theta) + C (a' (u), u,
   \Delta \theta) + R_a (u) \circ \Delta \theta . \]
Then, Ansatz (\ref{e:firstansatz}) gives
\[ a (u) \circ \Delta \theta = a' (u) (\theta \circ \Delta \theta) + a' (u)
   (u^{\sharp} \circ \Delta \theta) + C (a' (u), u, \Delta \theta) + R_a (u)
   \circ \Delta \theta . \]
Summarizing, we have:
\begin{eqnarray*}
  \Phi (u) & = & a' (u) (\theta \circ \Delta \theta) + a (u) \succ \Delta u +
  a' (u) (u^{\sharp} \circ \Delta \theta)\\
  &  & + C (a' (u), u, \Delta \theta) + R_a (u) \circ \Delta \theta + a (u)
  \circ \Delta u^{\sharp}
\end{eqnarray*}
Thanks to the nonlinear commutator (Lemma~\ref{l:nonlin-comm}), we can
decompose the resonant term $\theta \circ \Delta \theta$ to obtain
\begin{eqnarray*}
  \Phi (u) & = & a (u) \succ \Delta u + a' (u) (u^{\sharp} \circ \Delta
  \theta) + C (a' (u), u, \Delta \theta) + R_a (u) \circ \Delta \theta + a (u)
  \circ \Delta u^{\sharp}\\
  &  & + a' (u) \Lambda (a (u), \vartheta) + a' (u) \Pi_{\diamondsuit} (a
  (u), \Theta_2)
\end{eqnarray*}
and $\Lambda (a (u), \vartheta) \in \CC^{3 \alpha - 2 - \varepsilon}_T$ if $u
\in \LL^{\alpha}_T$. Here we defined
\begin{equation}
  \Theta_2 (\eta, x) \assign (\vartheta \circ \Delta \vartheta) (\eta, x) =
  \sum_{i \sim j} \Delta_i \vartheta (\eta, \cdummy) (x) \Delta_j [\Delta
  \vartheta (\eta, \cdummy)] (x) \label{e:def-Theta}
\end{equation}
Finally, recalling the decomposition of $u^{\sharp}$ in two
terms~(\ref{eq:decomp-sharp}) we obtain
\[ \Phi (u) = a' (u) \Pi_{\diamondsuit} (a (u), \Theta_2) + \Phi_1 (u) +
   \Phi_2 (u) \]
where
\begin{eqnarray*}
  \Phi_1 (u) & \assign & a (u) \succ \Delta u + C (a' (u), u, \Delta \theta) +
  R_a (u) \circ \Delta \theta + a' (u) \Lambda (a (u), \vartheta)\\
  &  & + a' (u) (\Pi_{\precprec} (a (u), U^{\sharp}) \circ \Delta \theta) + a
  (u) \circ \Delta \Pi_{\precprec} (a (u), U^{\sharp}),\\
  \Phi_2 (u) & \assign & a' (u) \left( \Pi_{\precprec} \left( a (u), \PP
  u_0^{\sharp} \right) \circ \Delta \theta \right) + a (u) \circ \Delta
  \Pi_{\precprec} \left( a (u), \PP u_0^{\sharp} \right) .
\end{eqnarray*}
Thanks to Lemma~\ref{l:decomp-est} the terms $a' (u) \Pi_{\diamondsuit} (a
(u), \Theta_2)$ and $\Phi_1 (u)$ can be estimated in $\CC^{2 \alpha - 2}_T$,
provided $\Theta_2 \in C^2_{\eta} \CC^{2 \alpha - 2}_T$ (see
Section~\ref{s:renorm}). On the other hand the term $\Phi_2 (u) (t)$ can be
estimated in $\CC^{2 \alpha - 2}$ only for strictly positive times $t > 0$ due
to the lack of regularity of the initial condition $u_0^{\sharp}$ which a
priori lives only in $\CC^{\alpha}$.

Note moreover that the specific form of $\Phi$ allows to deduce that if we
replace $\Theta_2$ by $\tilde{\Theta}_2 = \Theta_2 - H$ with $H \in C_{\eta}^2
\CC^{2 \alpha - 2}_T$ then this is equivalent to consider an equation for $u$
of the form
\[ \partial_t u (t, x) - a (u (t, x)) \Delta u (t, x) = \xi (x) - a' (u (t,
   x)) H (a (u (t, x)), t, x) . \]
Let us resume this long discussion in the following theorem:

\begin{theorem}
  \label{th:identification}Assume that $\xi \in \CC^0, u_0 \in \CC^2, H \in
  C_{\eta}^2 \CC_T^0$. $u \in C^1_T \CC^2$ is the classical solution to the
  equation
  \begin{equation}
    \partial_t u (t, x) - a (u (t, x)) \Delta u (t, x) = \xi (x) - a' (u (t,
    x)) H (a (u (t, x)), t, x), \qquad u (0) = u_0, \label{eq:pde}
  \end{equation}
  up to time $T > 0$ if
  \[ u = \Pi_{\precprec} \left( a (u), \vartheta + U^{\sharp} + \PP
     u_0^{\sharp} \right), \]
  where $\vartheta$ is the solution to eq.~(\ref{eq:def-theta}) and
  $U^{\sharp}$ is the solution to the PDE
  \begin{equation}
    (\partial_t - \eta \Delta) U^{\sharp} (\eta) = F (u, U^{\sharp},
    u_0^{\sharp}) \qquad U^{\sharp} (\eta, 0) = 0 \qquad \eta \in [\lambda, 1]
    \label{eq:pde-sharp}
  \end{equation}
  with
  \begin{eqnarray*}
    F (u, U^{\sharp}, u_0) & = & a' (u) \Pi_{\nocomma \diamondsuit} (a (u),
    \Theta_2) + \Phi_1 (u) + \Phi_2 (u) + \Psi (a (u), \vartheta) + \Psi (a
    (u), U^{\sharp})\\
    &  & + \Psi \left( a (u), \PP u_0^{\sharp} \right)
  \end{eqnarray*}
  and $\Theta_2 = \vartheta \circ \Delta \vartheta - H$.
\end{theorem}

\begin{definition}
  For any $\alpha \in \mathbbm{R}$ we define $\mathcal{X}^{\alpha} \subseteq
  C^2_{\eta} \CC^{\alpha} \times C^2_{\eta} \CC^{2 \alpha - 2}$ the closure of
  the image of the map
  \[ (\rho, H) \in C^2_{\eta} \CC^2 \times C^2_{\eta} \CC^0 \mapsto J (\rho,
     H) = (\rho, \rho \circ \Delta \rho - H) \in C^2_{\eta} \CC^2 \times
     C^2_{\eta} \CC^0 \]
  (in the topology of $C^2_{\eta} \CC^{\alpha} \times C^2_{\eta} \CC^{2 \alpha
  - 2}$).
\end{definition}

We call the elements in $\mathcal{X}^{\alpha}$ \tmtextit{enhanced noises}. In
the next section we will exploit the space $\mathcal{X}^{\alpha}$ for $2 / 3 <
\alpha < 1$ to solve equations~(\ref{eq:pde-sharp}) and (\ref{e:firstansatz}).
\

\section{Local wellposedness\label{s:fixpoint}}

The main result of this section is the local well--posedness for
equations~(\ref{e:firstansatz}) and~(\ref{eq:pde-sharp}) when $(\vartheta,
\Theta_2) \in \mathcal{X}^{\alpha}$ and $u_0 \in \CC^{\alpha}$ for $2 / 3 <
\alpha < 1$. This yields a unique solution to~(\ref{eq:pde}), thanks to
Theorem~\ref{th:identification}.

\begin{theorem}
  Let $\alpha > 2 / 3.$ Then for any $(\vartheta, \Theta_2) \in
  \mathcal{X}^{\alpha}$ and $u_0 \in \CC^{\alpha}$ there exists a time $T > 0$
  depending only on $\| (\vartheta, \Theta_2) \|_{\mathcal{X}^{\alpha}}$ and
  $\| u_0 \|_{\alpha}$ up to which the system of
  equations~(\ref{e:firstansatz}) and~(\ref{eq:pde-sharp}) has a unique
  solution $(u, U^{\sharp}) \in \LL^{\alpha}_T \times C_{\eta}^2 \LL^{2
  \delta}_T$ for all $\delta < \alpha$ such that $2 \delta + \alpha > 2$. For
  any fixed $\tau > 0$ there exist a ball $B_{\tau} \subseteq \CC^{\alpha}
  \times \mathcal{X}^{\alpha}$ such that the solution map
  \[ \Sigma_{\tau} : (u_0, \vartheta, \Theta_2) \in B_{\tau} \mapsto (u,
     U^{\sharp}) \in \LL^{\alpha}_{\tau} \times C_{\eta}^2 \LL^{2
     \delta}_{\tau} \]
  is well defined and Lipshitz continuous in the data. 
\end{theorem}

\begin{proof}
  Let $\mathcal{G}_T = \LL^{\alpha}_T \times C_{\eta}^2 \LL^{2 \delta}_T$. We
  introduce the map
  \[ \Gamma : (u, U^{\sharp}) \in \mathcal{G}_T \mapsto (\Gamma_u (u,
     U^{\sharp}), \Gamma_{U^{\sharp}} (u, U^{\sharp})) \in \mathcal{G}_T \]
  by
  \[ \Gamma_u (u, U^{\sharp}) \assign \Pi_{\precprec} (a (u), \vartheta) +
     \Pi_{\precprec} (a (u), \Gamma_{U^{\sharp}} (u, U^{\sharp})) +
     \Pi_{\precprec} \left( a (u), \PP u_0^{\sharp} \right) \]
  and
  \[ (\partial_t - \eta \Delta) \Gamma_{U^{\sharp}} (u, U^{\sharp}) (\eta) = F
     (u, U^{\sharp}, u_0^{\sharp}), \qquad \Gamma_{U^{\sharp}} (u, U^{\sharp})
     (\eta) (0) = 0, \qquad \eta \in [\lambda, 1], \]
  We will establish that this map is a contraction in the space $\mathcal{G}_T
  .$
  
  First, we have to show that there exists a ball $B \subset \mathcal{G}_T$
  such that $\Gamma (B) \subseteq B$. We have the bound $\left\| \PP
  u_0^{\sharp} \right\|_{C^1_{\eta} \LL^{\alpha}_T} \lesssim \| u^{\sharp}_0
  \|_{\CC^{\alpha}_T}$ . It is easy to obtain, using the estimates of
  Section~\ref{s:nonlinear} and Lemma~\ref{l:schauder}:
  \begin{eqnarray*}
    &  & \left\| \int_0^T e^{- \eta \Delta (t - s)} \left[ \Phi_1 (u) + \Psi
    (a (u), \vartheta) + \Psi (a (u), U^{\sharp}) + \Psi \left( a (u), \PP
    u_0^{\sharp} \right) \right]_s \right\|_{C^2_{\eta} \LL_T^{2 \delta}}\\
    & \lesssim & T^{\kappa} \left( 1 + \| u \|_{\LL^{\alpha}_T} \right)^4
    \left( 1 + \| \xi \|_{\CC^{\alpha - 2}} \right)^2 \| u_0^{\sharp}
    \|_{\CC^{\alpha}} \left( 1 + \| U^{\sharp} \|_{C^2_{\eta} \LL^{2
    \delta}_T} \right)
  \end{eqnarray*}
  for some $\kappa > 0$.
  
  By the assumption that $(\vartheta, \Theta_2) \in \mathcal{X}^{\alpha}$ we
  deduce that there exists $M > 0$ such that $\| \Theta_2 \|_{C_{\eta}^2
  \CC_T^{2 \alpha - 2}} \leqslant M$. We have
  \[ \left\| \int_0^T e^{- \eta \Delta (t - s)} [a' (u) \Pi_{\nocomma
     \diamondsuit} (a (u), \Theta_2)]_s \right\|_{C^2_{\eta} \LL_T^{2 \delta}}
     \lesssim T^{\alpha - \delta} \left( 1 + \| u \|_{\CC^{\alpha}_T}
     \right)^2 \| \Theta_2 \|_{C_{\eta}^2 \CC_T^{2 \alpha - 2}} . \]
  To bound the term $\Phi_2 (u)$ we observe that $\left\| \PP_t u^{\sharp}_0
  \right\|_{C_{\eta}^2 \CC^{2 \alpha}} \lesssim t^{- \frac{\alpha}{2}} \|
  u^{\sharp}_0 \|_{\CC^{\alpha}}$ thanks to Lemma~\ref{l:schauder}. This gives
  \[ \left\| \int_0^T e^{- \eta \Delta (t - s)} \Phi_2 (u)_s
     \right\|_{C^2_{\eta} \LL_T^{2 \delta}} \lesssim T^{\alpha - \delta}
     \left( 1 + \| u \|_{\CC^{\alpha}_T} \right) \left( 1 + \| \xi
     \|_{\CC^{\alpha - 2}} \right) \| u^{\sharp}_0 \|_{\CC^{\alpha}} \]
  and then $\Gamma_{U^{\sharp}} (u, U^{\sharp})$ is bounded in $C^2_{\eta}
  \LL_T^{2 \delta}$ for $T$ small enough. We have also
  \begin{eqnarray*}
    \| \Gamma_u (u, U^{\sharp}) \|_{\LL^{\alpha}_T} & \lesssim & \| \xi
    \|_{\CC^{\alpha - 2}} + \| u^{\sharp}_0 \|_{\CC^{\alpha}} + \|
    \Gamma_{U^{\sharp}} (u, U^{\sharp}) \|_{C_{\eta} \LL_T^{\alpha}}\\
    & \lesssim & \| \xi \|_{\CC^{\alpha - 2}} + \| u^{\sharp}_0
    \|_{\CC^{\alpha}} + T^{\frac{2 \delta - \alpha}{2}} \| \Gamma_{U^{\sharp}}
    (u, U^{\sharp}) \|_{C^2_{\eta} \LL_T^{2 \delta}}
  \end{eqnarray*}
  and these bounds show that $\Gamma (B) \subseteq B$. The contractivity of
  $\Gamma_{U^{\sharp}} (u, U^{\sharp})$ can be obtained in the same way. Now
  consider $\Gamma_u (u, U^{\sharp})$: we have
  \begin{equation*}
  | \Pi_{\precprec} (a (u_1), U_1^{\sharp}) - \Pi_{\precprec} (a
    (u_2), U_2^{\sharp}) \|_{\LL^{\alpha}_T}
     \lesssim  T^{\frac{2 \delta - \alpha}{2}} \left( \| U^{\sharp}_1 -
    U^{\sharp}_2 \|_{C_{\eta} \LL^{2 \delta}} + \| u_1 - u_2 \|_{C_T
    L^{\infty}} \| U^{\sharp}_2 \|_{C_{\eta}^1 \LL^{2 \delta}} \right)
  \end{equation*}
%  \begin{eqnarray*}
%    &  & \| \Pi_{\precprec} (a (u_1), U_1^{\sharp}) - \Pi_{\precprec} (a
%    (u_2), U_2^{\sharp}) \|_{\LL^{\alpha}_T}\\
%    & \lesssim & T^{\frac{2 \delta - \alpha}{2}} \left( \| U^{\sharp}_1 -
%    U^{\sharp}_2 \|_{C_{\eta} \LL^{2 \delta}} + \| u_1 - u_2 \|_{C_T
%    L^{\infty}} \| U^{\sharp}_2 \|_{C_{\eta}^1 \LL^{2 \delta}} \right)
%  \end{eqnarray*}
  while for the other terms in $\Gamma_u (u_1, U^{\sharp}_1) - \Gamma_u (u_2,
  U^{\sharp}_2)$ we remark that
  \[ \sup_{s \in [0, t]} \| u_{1, s} - u_0 - u_{2, s} + u_0 \|_{L^{\infty}}
     \lesssim t^{\varepsilon / 2} \| u_1 - u_2 \|_{C_{[0, t]}^{\varepsilon /
     2} L^{\infty}} . \]
  Then $\forall 0 < \varepsilon < \alpha$, using Lemma~\ref{l:nonlinear-est}
  and Lemma~\ref{l:interpol}:
  \begin{eqnarray*}
    \| \Pi_{\precprec} (a (u_1), \vartheta) - \Pi_{\precprec} (a (u_2),
    \vartheta) \|_{\LL^{\alpha}_T} & \lesssim & \| a (u_1) - a (u_2) \|_{C_T
    L^{\infty}} \| \vartheta \|_{C_{\eta}^1 \LL^{\alpha}_T}\\
    & \lesssim & \| u_1 - u_2 \|_{C_T L^{\infty}} \| \xi \|_{\CC^{\alpha -
    2}}\\
    & \lesssim & T^{\varepsilon / 2}  \| u_1 - u_2 \|_{C_T^{\varepsilon / 2}
    L^{\infty}} \| \xi \|_{\CC^{\alpha - 2}} \\
    & \lesssim & T^{\varepsilon / 2} \| u_1 - u_2 \|_{\LL^{\alpha}_T} \| \xi
    \|_{\CC^{\alpha - 2}} .
  \end{eqnarray*}
  With the same reasoning we estimate
  \begin{eqnarray*}
    \left\| \Pi_{\precprec} \left( a (u_1), \PP u^{\sharp}_0 \right) -
    \Pi_{\precprec} \left( a (u_2), \PP u^{\sharp}_0 \right)
    \right\|_{\LL^{\alpha}_T} & \lesssim & T^{\varepsilon / 2} \| u_1 - u_2
    \|_{C^{\varepsilon / 2}_T L^{\infty}} \left\| \PP u^{\sharp}_0
    \right\|_{C_{\eta}^1 \LL^{\alpha}_T}\\
    & \lesssim & T^{\varepsilon / 2} \| u_1 - u_2 \|_{\LL^{\alpha}_T} \|
    u^{\sharp}_0 \|_{\CC^{\alpha}_T}
  \end{eqnarray*}
  and then $\Gamma$ is a contraction for small times.
  
  The uniqueness of the solution $(u, U^{\sharp}) \in \LL^{\alpha}_T \times
  C_{\eta}^2 \LL^{2 \delta}_T$ and the Lipshitz continuity of the localized
  solution map $\Sigma_{\tau}$ can be proved along the same lines via standard
  arguments.
\end{proof}

\section{Renormalization\label{s:renorm}}

At this point we want to construct \tmtextit{an} enhanced noise $\Xi$
associated to the white noise $\xi$. Already in the standard setting of the
generalised PAM model with constant diffusion matrix, the construction of the
enhancement requires a renormalization since the resonant product $\vartheta
\circ \Delta \vartheta$ is not well defined.

Let $\psi \in \mathcal{S} (\mathbbm{T}^2)$ be a cutoff function and let
$\psi_{\varepsilon} (x) = \varepsilon^{- 2} \psi (x / \varepsilon)$. Then
define a regularised noise by $\xi_{\varepsilon} = \psi_{\varepsilon} \ast
\xi$ and let $\vartheta_{\varepsilon} = (- \eta \Delta)^{- 1}
\xi_{\varepsilon}$. Notice that
\[ H_{\varepsilon} (\eta) \assign \mathbbm{E} [\vartheta_{\varepsilon} (\eta,
   x) \circ \Delta \vartheta_{\varepsilon} (\eta, x)] =\mathbbm{E}
   [\vartheta_{\varepsilon} (\eta, x) \Delta \vartheta_{\varepsilon} (\eta,
   x)] = - \sum_{k \in \mathbbm{Z}^2 \backslash \{ 0 \}}
   \frac{\hat{\psi}_{\varepsilon} (k)^2}{\eta^2 k^2} = -
   \frac{\sigma_{\varepsilon}}{\eta^2} \]
where
\[ \sigma_{\varepsilon} \assign \sum_{k \in \mathbbm{Z}^2 \backslash \{ 0 \}}
   \frac{\hat{\psi}_{\varepsilon} (k)^2}{k^2} \simeq | \log \varepsilon | \]
as $\varepsilon \rightarrow 0$. Subtracting the diverging quantity
$H_{\varepsilon}$ to $\vartheta_{\varepsilon} \circ \Delta
\vartheta_{\varepsilon}$ and then taking the limit as $\varepsilon \rightarrow
0$ delivers a finite result.

\begin{theorem}
  \label{th:conv-white-noise}Take $\alpha < 1$ and let $\Xi_{\varepsilon} =
  (\xi_{\varepsilon}, \Xi_{2, \varepsilon}) : = (\xi_{\varepsilon},
  \vartheta_{\varepsilon} \circ \Delta \vartheta_{\varepsilon} -
  H_{\varepsilon})$. Then the family $(\Xi_{\varepsilon})_{\varepsilon}
  \subseteq \mathcal{X}^{\alpha}$ converges a.s. and in $L^p$ to a random
  element $\Xi = (\xi, \Xi_2) \in \mathcal{X}^{\alpha}$.
\end{theorem}

\begin{proof}
  The proof is a mild modification of the proof for
  PAM~{\cite{gubinelli_paracontrolled_2015}}. In order to establish the
  required $C^2_{\eta} \CC^{2 \alpha - 2}_T$ regularity for $\Xi_2$ we follow
  the computations for the case where the diffusion coefficient is constant.
  We have only to discuss the additional regularity in the $\eta$ parameter.
  In order to do so observe that
  \[ \Xi_{2, \varepsilon} (\eta) = \sum_{i \sim j} \llbracket \Delta_i
     \vartheta_{\varepsilon} (\eta) \Delta_j \Delta \vartheta_{\varepsilon}
     (\eta) \rrbracket \]
  where $\llbracket \rrbracket$ denotes the Wick product with respect to the
  Gaussian structure of $\xi$. Then we have
  \[ \partial_{\eta} \Xi_{2, \varepsilon} (\eta) = \sum_{i \sim j} \llbracket
     \Delta_i \partial_{\eta} \vartheta_{\varepsilon} (\eta) \Delta_j \Delta
     \vartheta_{\varepsilon} (\eta) \rrbracket + \sum_{i \sim j} \llbracket
     \Delta_i \vartheta_{\varepsilon} (\eta) \Delta_j \Delta \partial_{\eta}
     \vartheta_{\varepsilon} (\eta) \rrbracket, \]
  and
    \begin{eqnarray*}
    \partial_{\eta}^2 \Xi_{2, \varepsilon} (\eta) & = & \sum_{i \sim j}
    \llbracket \Delta_i \partial_{\eta}^2 \vartheta_{\varepsilon} (\eta)
    \Delta_j \Delta \vartheta_{\varepsilon} (\eta) \rrbracket + \sum_{i \sim
    j} \llbracket \Delta_i \vartheta_{\varepsilon} (\eta) \Delta_j \Delta
    \partial_{\eta}^2 \vartheta_{\varepsilon} (\eta) \rrbracket\\
    &  & + \sum_{i \sim j} 2 \llbracket \Delta_i \partial_{\eta}
    \vartheta_{\varepsilon} (\eta) \Delta_j \Delta \partial_{\eta}
    \vartheta_{\varepsilon} (\eta) \rrbracket .
  \end{eqnarray*}
  
  
  Now the computations relative to the regularities of these additional
  stochastic objects are equivalent to those for the term $\Xi_{2,
  \varepsilon}$ where one or two instances of $\vartheta_{\varepsilon} (\eta)$
  are replaced by Gaussian fields of similar regularities of the form
  $\partial_{\eta} \vartheta_{\varepsilon} (\eta)$ and $\partial_{\eta}^2
  \vartheta_{\varepsilon} (\eta)$, a direct inspection of the proof allows us
  to deduce that we have almost sure $\CC^{2 \alpha - 2}$ regularity for these
  terms and also for random fields $\partial_{\eta}^n \Xi_{2, \varepsilon}$
  for any finite $n$. This allows also to deduce that the random field is a.s.
  smooth in the parameter $\eta$. Similar computations allow to prove
  continuity in $\varepsilon$ for $\varepsilon > 0$. The rest of the proof is
  standard.
\end{proof}

In conclusion we see that in order to be able to use this convergence result
we need to modify our approximate PDE and consider instead
\[ \partial_t u_{\varepsilon} - a (u_{\varepsilon}) \Delta u_{\varepsilon} =
   \xi_{\varepsilon} - a' (u_{\varepsilon}) H_{\varepsilon} (a
   (u_{\varepsilon})) \]
which gives the renormalised equation (\ref{eq:ur-formulation-approx}).

Our well--posedness results for the paracontrolled formulation of this
equation together with the convergence result in
Theorem~\ref{th:conv-white-noise} allow to deduce that $u_{\varepsilon}
\rightarrow u$ in $\CC^{\delta}_T$ for any $2 / 3 < \delta < \alpha < 1$ and
that the limiting process $u$ satisfies a modified version of
eq.~(\ref{eq:ur-formulation}), namely
\[ \partial_t u - a (u) \diamond \Delta u = \xi, \qquad u (0) = u_0, \]
where $a (u) \diamond \Delta u$ denotes a \tmtextit{renormalized} diffusion
term given by
\begin{equation}
  a (u) \diamond \Delta u \assign a (u) \prec \Delta u + a' (u)
  \Pi_{\diamondsuit} (a (u), \Xi_2) + \Phi_1 (u) + \Phi_2 (u) .
  \label{eq:ren-diff}
\end{equation}
\section{Nonlinear source terms}\label{sec:gen-nonlinear}

Let us start by discussing the presence of a $u$ dependent r.h.s. in
eq.~(\ref{eq:ur-formulation}). We want to solve
\[ \partial_t u - a_1 (u) \Delta u = a_2 (u) \xi \]
where $a_1$ is a non-linear diffusion coefficient as before and $a_2 :
\mathbbm{R} \rightarrow \mathbbm{R}$ is another bounded function with
sufficiently many bounded derivatives. We rewrite this equation as
\[ \Pi_{\prec} \left( a (u), \LL \right) u = a_2 (u) \prec \xi + a_1 (u) \circ
   \Delta u + a_2 (u) \circ \xi + a_1 (u) \succ \Delta u + a_2 (u) \succ \xi
\]
where now $a (u) = (a_1 (u), a_2 (u))$ is a vector valued non-linearity. Since
we don't need $u$ to depend on any parameter $\eta = (\eta_1, \eta_2)$, we
have defined $\LL$ as
\[ \LL (\eta) \assign \partial_t - \eta_1 \Delta \]
and used the identity $\Pi_{\prec} \left( a (u), \LL \right) u = (\partial_t -
a_1 (u) \prec \Delta) u$, similarly to what we have done in
(\ref{e:paradiff-oper}).

Notice that the non--linear paraproduct can be extended trivially to the
vector valued case in such a way that, for example,
\[ \Pi_{\precprec} ((g_1, g_2), h) (t, x) = \sum_i \int_{y, s} Q_{i, t} (s)
   P_{i, x} (y) (\Delta_i h ((g_1 (s, y), g_2 (s, y)), t, \cdot)) (x) . \]
As before we make the Ansatz
\[ u = \Pi_{\precprec} (a (u), \vartheta) + u^{\sharp} \]
where now $\vartheta$ solves
\[ \LL (\eta) \vartheta (\eta, x) = (\partial_t - \eta_1 \Delta) \vartheta
   (\eta, x) = \eta_2 \xi (x), \]
for $\eta = (\eta_1, \eta_2) \in [\lambda, 1] \times [- L, L]$ where $L$ is a
large but fixed constant. The bounded domain is important to be able to have
uniform estimates and reuse the estimates proved above in the simple situation
of $\eta_2 = 1$. The solution of this equation is
\[ \vartheta (\eta, \cdot) = \eta_2 \int_0^{\infty} e^{\eta_1 \Delta s} \xi
   \mathd s = - \frac{\eta_2}{\eta_1} \Delta^{- 1} \xi . \]
Observe that
$$
\Pi_{\prec} \left( a (u), \LL \right) u = \Pi_{\prec} \left( a (u), \LL
   \right) \Pi_{\precprec} (a (u), \vartheta) + \Pi_{\prec} \left( a (u), \LL
   \right) u^{\sharp}
$$
and recall that (Lemma~\ref{l:def-psi}) $\Pi_{\prec} \left( a (u), \LL \right)
\Pi_{\precprec} (a (u), \vartheta) = \Pi_{\precprec} \left( a (u), \LL
\vartheta \right) + \Psi (a (u), \vartheta)$. Now
$$
\left( \LL \vartheta \right) (\eta) = (\partial_t - \eta_1 \Delta)
   \vartheta (\eta, t, x) = \Xi (\eta), \qquad \eta = (\eta_1, \eta_2) \in
   [\lambda, 1] \times [- L, L]
$$
with $\Xi (\eta) (t, x) = \eta_2 \xi (x)$ and then
\[ \Pi_{\precprec} \left( a (u), \LL \vartheta \right) = \Pi_{\precprec} (a
   (u), \Xi) = a_2 (u) \precprec \xi . \]
In conclusion
\[ \Pi_{\prec} \left( a (u), \LL \right) \Pi_{\precprec} (a (u), \vartheta) =
   a_2 (u) \precprec \xi + \Psi (a (u), \vartheta) \]
and the equation for $u^{\sharp}$ reads
\begin{eqnarray}
  \Pi_{\prec} \left( a (u), \LL u^{\sharp} \right) & = & a_1 (u) \circ \Delta
  u + a_2 (u) \circ \xi + [a_2 (u) \prec \xi - a_2 (u) \precprec \xi]
  \nonumber\\
  &  & + a_1 (u) \succ \Delta u + a_2 (u) \succ \xi - \Psi (a (u), \vartheta)
  \nonumber
\end{eqnarray}
where now all the terms on the r.h.s. can be treated as remainder terms. Let
us just remark that the commutation term $a_2 (u) \prec \xi - a_2 (u)
\precprec \xi$ has a standard treatment via Lemma~\ref{l:comm-paraproduct}. Of
course, the first two terms require to be treated as resonant terms. Note
that, modulo terms of order $\CC^{3 \alpha - 2}_T$ (or $\mathcal{E}^{\alpha /
2} \CC^{2 \alpha - 2}$) the terms $a_1 (u) \circ \Delta u + a_2 (u) \circ \xi$
are equivalent to
\[ a'_1 (u) \Pi_{\diamondsuit} (a (u), \vartheta \circ \Delta \vartheta) +
   a'_2 (u) (\Pi_{\precprec} (a (u), \vartheta) \circ \xi) \]
and that by computations similar to those of the previous sections one can
prove that
\[ (\Pi_{\precprec} (a (u), \vartheta) \circ \xi) = \Pi_{\diamondsuit} (a (u),
   \vartheta \circ \xi) + \CC^{3 \alpha - 2}_T \]
so the resonant terms are comparable to the sum of the two terms
\[ a'_1 (u) \Pi_{\diamondsuit} (a (u), \vartheta \circ \Delta \vartheta) +
   a'_2 (u) \Pi_{\diamondsuit} (a (u), \vartheta \circ \xi) \]
which require renormalization of the form
\begin{equation}
  \frac{a'_1 (u) a_2 (u)^2}{a_1 (u)^2} \sigma_{\varepsilon} - \frac{a'_2 (u)
  a_2 (u)}{a_1 (u)} \sigma_{\varepsilon} \label{eq:renorm-non-linear-rhs}
\end{equation}
and the convergence follows with the same arguments of Section~\ref{s:renorm}.

We remark that the structure of the second renormalisation term, which is due
to the r.h.s. in the equation, is the same of that found by Bailleul,
Debussche and Hofmanova in~{\cite{hofmanova}}.

\begin{remark}
  \label{r:vector-form}Our approach works straightforwardly for the equation
  \[ \partial_t u (t, x) - a_{i j} (u (t, x)) \partial^2_{i j} u (t, x) = g (u
     (t, x)) \xi \]
  with $a : \mathbbm{R} \rightarrow M_2 (\mathbbm{R})$ such that $\sum_{i, j}
  a (u)_{i j} x_i x_j \geqslant C | x |^2$ $\forall x \in \mathbbm{R}^2$ for
  $C > 0$ and $\partial^2_{i j} \assign \frac{\partial^2}{\partial x_i
  \partial x_j}$.
  
  To see that, let $a (u) \assign (a_{i j} (u), g (u)) \in \mathbbm{R}^5$ and
  $\eta = (\eta_{i, j}, \eta_g) \in \mathbbm{R}^5$. Let $\LL (\eta) \assign
  \partial_t - \eta_{i j} \partial^2_{i j}$ and $\Xi (\eta) \assign \eta_g
  \xi$ with the uniform ellipticity condition $\sum_{i, j} \eta_{i j} x_i x_j
  \geqslant C | x |^2$ $\forall x \in \mathbbm{R}^2$. It is easy to verify
  that Lemma~\ref{l:def-psi} and Lemma~\ref{l:nonlin-comm} hold within this
  setting, just considering nonlinear paraproducts for functions depending on
  5 parameters. We have then:
  \[ u = \Pi_{\precprec} \left( a (u), \vartheta + U^{\sharp} + \PP
     u_0^{\sharp} \right) \]
  with $\vartheta (\eta)$ stationary solution of $\LL \vartheta (\eta) = \Xi
  (\eta)$, $\PP_t u_0^{\sharp} \assign e^{\eta_{i j} \partial^2_{i j} t}
  u_0^{\sharp}$ and $U^{\sharp} (\eta)$ which solves
  \[ \LL U^{\sharp} (\eta) = \Pi_{\diamondsuit} ((a (u), a' (u)), \Theta_1) +
     \Pi_{\diamondsuit} ((a (u), a' (u)), \Theta_2) + Q (u, U^{\sharp}) \]
  with $Q (u, U^{\sharp}) \in \CC^{2 \alpha - 2 - \varepsilon}$, $\Theta_1 \in
  C^2_{\eta} C^2_{\eta'} \CC^{2 \alpha - 2} = \vartheta (\eta) \circ \eta'_{i
  j} \partial^2_{i j} \vartheta (\eta)$, $\Theta_2 (\eta, \eta') = \vartheta
  (\eta) \circ \eta'_g \xi$ and $U^{\sharp} (t = 0) = 0$. Note that we can
  write $\vartheta$ as
  \[ \vartheta (\eta) = \eta_g \int_0^{\infty} e^{t \eta_{i j} \partial^2_{i
     j}} \xi \mathd t \qquad \hat{\vartheta} (k) = \eta_g \frac{\hat{\xi}
     (k)}{\eta_{i j} k_i k_j}, \quad k \in \mathbbm{Z}^2 \backslash \{ 0 \} .
  \]
  From the uniform ellipticity condition we have that $\| \vartheta
  \|_{C_{\eta}^k \CC^{\alpha}} \lesssim \| \xi \|_{\CC^{\alpha}}$, and
  Schauder estimates analogous to those of Lemma~\ref{l:schauder} hold as
  well.
  
  Now consider the renormalization. We have
  \begin{eqnarray*}
    H_1^{\varepsilon} (\eta, \eta') \assign \mathbbm{E} (\Theta_1 (\eta,
    \eta')) & = & - \eta_g^2 \sum_{k \in \mathbbm{Z}^2 \backslash \{ 0 \}}
    \hat{\psi}_{\varepsilon} (k)^2 \frac{\sum_{i, j} \eta'_{i j} k_i
    k_j}{\left( \sum_{i, j} \eta_{i j} k_i k_j \right)^2},\\
    H_2^{\varepsilon} (\eta, \eta') \assign \mathbbm{E} (\Theta_2 (\eta,
    \eta')) & = & \eta_g \eta'_g \sum_{k \in \mathbbm{Z}^2 \backslash \{ 0 \}}
    \frac{\hat{\psi}_{\varepsilon} (k)^2}{\sum_{i, j} \eta_{i j} k_i k_j} .
  \end{eqnarray*}
  The convergence of $\Theta_1^{\varepsilon} - H_1^{\varepsilon}$,
  $\Theta_2^{\varepsilon} - H_2^{\varepsilon}$ in $C_{(\eta, \eta')}^k \CC^{2
  \alpha - 2} (\mathbbm{T}^2)$ can be obtained with the techniques used in
  {\cite{gubinelli_paracontrolled_2015}}, Section 5.2. 
\end{remark}

\section{Full generality}

Within the framework of the present work we are actually able to treat
equations of the form
\begin{equation}
  \partial_t u (t, x) - a_1 (u (t, x)) \Delta u (t, x) = \xi (a_2 (u (t, x)),
  x) \label{eq:general}
\end{equation}
where $\xi (\eta_2, x)$ is a Gaussian process with covariance
\[ \mathbbm{E} [\xi (\eta_2, x) \xi (\tilde{\eta}_2, \tilde{x})] = F (\eta_2, \tilde{\eta}_2)
   \delta (x - \tilde{x}) \]
where $F$ is a smooth covariance function. Let as before $2 / 3 < \alpha < 1$.
In this case we can take as skeleton equation
\[ \LL (\eta) \vartheta \assign \partial_t \vartheta (\eta, t, x) - \eta_1
   \Delta \vartheta (\eta, t, x) = \xi (\eta_2, x) \]
whose solution $\vartheta$ is a Gaussian process smooth in the variable $\eta
= (\eta_1, \eta_2)$ which we assume taking value in a compact subset of
$\mathbbm{R}^2$ for which $\eta_1 \geqslant \lambda > 0$ with fixed $\lambda$.
Letting $a (u) = (a_1 (u), a_2 (u))$ we can rewrite the l.h.s. of
eq.~(\ref{eq:general}) in the form
\[ \partial_t u - a_1 (u) \Delta u = \Pi_{\diamondsuit} \left( a (u), \LL u
   \right) \]
and the r.h.s. as
\[ \xi (a_2 (u (t, x)), x) = \Pi_{\diamondsuit} (a (u), \Xi) \]
where $\Xi (\eta, x) = \xi (\eta_2, x)$. Now we perform the paraproduct
decomposition to get
\begin{eqnarray*}
  \Pi_{\prec} \left( a (u), \LL u \right) - \Pi_{\prec} (a (u), \Xi) & = &
  \Pi_{\circ} (a (u), \Xi) + \Pi_{\circ} \left( a (u), \DF u \right)\\
  &  & + \Pi_{\succ} (a (u), \Xi) + \Pi_{\succ} \left( a (u), \DF u \right) .
\end{eqnarray*}
with $\DF (\eta) \assign \eta_1 \Delta$. Let $\PP_t (\eta) \assign e^{t \eta_1
\Delta}$ as before, and invoke the paracontrolled Ansatz in the usual form
\[ u = \Pi_{\precprec} \left( a (u), \vartheta + U^{\sharp} + \PP
   u^{\sharp}_0 \right) . \]
Using that
\begin{eqnarray*}
  \Pi_{\prec} \left( a (u), \LL \Pi_{\precprec} \left( a (u), \vartheta +
  U^{\sharp} + \PP u^{\sharp}_0 \right) \right) & = & \Pi_{\precprec} \left( a
  (u), \LL \left( \vartheta + U^{\sharp} + \PP u_0^{\sharp} \right) \right)\\
  &  & + \Psi \left( a (u), \vartheta + U^{\sharp} + \PP u^{\sharp}_0 \right)
\end{eqnarray*}
and observing that we can take $\LL \vartheta = \Xi$ and that $\LL \PP
u^{\sharp}_0 = 0$ to get
\[ \Pi_{\precprec} \left( a (u), \LL U^{\sharp} \right) = F (u, U^{\sharp}) \]
where
\begin{eqnarray*}
  F (u, U^{\sharp}) & = & \Pi_{\circ} (a (u), \Xi) + \Pi_{\circ} \left( a (u),
  \DF u \right) + \Pi_{\succ} (a (u), \Xi) + \Pi_{\succ} \left( a (u), \DF u
  \right)\\
  &  & + [\Pi_{\prec} (a (u), \Xi) - \Pi_{\precprec} (a (u), \Xi)] - \Psi
  \left( a (u), \vartheta + U^{\sharp} + \PP u^{\sharp}_0 \right)
\end{eqnarray*}
which is solved by $U^{\sharp}$ satisfying
\[ \LL U^{\sharp} = F (u, U^{\sharp}) . \]
Indeed $\Pi_{\prec} (a (u), F (u, U^{\sharp})) = F (u, U^{\sharp})$, since $F
(u, U^{\sharp})$ does not depend on the additional parameter. Remark that the
term $\Pi_{\prec} (a (u), \Xi) - \Pi_{\precprec} (a (u), \Xi)$, which does not
appear in the simpler case, can be treated with
Lemma~\ref{l:comm-nonlin-para}.

It remains now to discuss the handling of the resonant products under the
paracontrolled assumption, namely $\Pi_{\circ} (a (u), \Xi)$ and $\Pi_{\circ}
\left( a (u), \DF u \right)$. Next lemma is a paralinearization result adapted
to our non-linear context.

\begin{lemma}
  Assume that $u \in \CC_T^{\rho}$ and $Z \in C_{\eta}^2 \CC^{\gamma}_T$ then
  if $\gamma + 2 \rho > 0$ we have
  \[ C (u, Z) \assign \Pi_{\circ} (a (u), Z) - u \circ \Pi_{\prec} ((a (u), a'
     (u)), \mathcal{D}Z) \in \CC_T^{\gamma + 2 \rho} \]
  where $\mathcal{D}_{} Z ((\eta, \eta'), t, x) \assign \sum_i \eta'_i
  \partial_{\eta_i} Z (\eta, t, x)$.
\end{lemma}

\begin{proof}
  
  \begin{eqnarray*}
    \Pi_{\circ} (a (u), Z) (t, x) & = & \sum_{i \sim j} \int_{y, z} K_{i, x}
    (y) K_{j, x} (z) Z (a (u (t, y)), t, z)\\
    & = & \sum_{i \sim j} \sum_k \int_{y, z, z', z''} K_{i, x} (y) K_{j, x}
    (z) P_{k, z} (z'') K_{k, z} (z') Z (a (u (t, y)), t, z')
  \end{eqnarray*}
  
\begin{eqnarray*}
  & = & \sum_{i \sim j} \sum_k \int_{y, z, z', z''} K_{i, x} (y) K_{j, x} (z)
     P_{k, z} (z'') K_{k, z} (z') [Z (a (u (t, y)), t, z') - Z (a (u (t,
     z'')), t, z')] \\
  & = & \sum_{i \sim j} \sum_k \int_{y, z, z', z''} K_{i, x} (y) K_{j, x} (z)
     P_{k, z} (z'') K_{k, z} (z')  \left[ \sum_{\ell} a'_{\ell} (u (t, z'')) \delta u^{t y}_{t z''} \partial_{a_{\ell}} Z (a (u (t, z'')), t, z') \right] \\
  & & + \sum_{i \sim j} \sum_{k \sim j} \int_{\underset{y, z}{z', z''}} K_{i,
     x} (y) K_{j, x} (z) P_{k, z} (z'') K_{k, z} (z') O ((\delta u^{t y}_{t z''})^2) \partial_{\eta}^2 Z (a (u (t, y)), t, z') 
\end{eqnarray*}     

  and observe that the first term is equal to $u \circ \Pi_{\prec} ((a (u), a'
  (u)), \mathcal{D}Z)$ while the second term can be easily estimated in
  $\CC_T^{\gamma + 2 \rho}$. 
\end{proof}

Using this result and Lemma~\ref{l:comm-nonlin-para} we can expand
\begin{eqnarray*}
  \Pi_{\circ} (a (u), \Xi) & = & u \circ \Pi_{\prec} ((a (u), a' (u)),
  \mathcal{D} \Xi) + \CC_T^{3 \alpha - 2}\\
  & = & \Pi_{\precprec} (a (u), \vartheta) \circ \Pi_{\prec} ((a (u), a'
  (u)), \mathcal{D} \Xi) + \CC_T^{3 \alpha - 2}\\
  & = & \Pi_{\diamondsuit} ((a (u), a' (u)), \vartheta \circ \mathcal{D} \Xi)
  + \CC_T^{3 \alpha - 2}
\end{eqnarray*}
and similarly, noting that
\begin{eqnarray*}
  \Pi_{\prec} \left( (a (u), a' (u)), \left( \mathcal{D} \DF \right) u \right)
  & = & \Pi_{\prec} \left( (a (u), a' (u)), \left( \mathcal{D} \DF \right)
  \Pi_{\precprec} (a (u), \vartheta) \right) + \CC_T^{3 \alpha - 2}\\
  & = & \Pi_{\prec} \left( (a (u), a' (u)), \left( \mathcal{D} \DF \right)
  \vartheta \right) + \CC_T^{3 \alpha - 2}
\end{eqnarray*}
where $\left( \mathcal{D} \DF \right) (\eta, \eta') = \eta_1' \Delta$, we have
\begin{eqnarray*}
  \Pi_{\circ} \left( a (u), \LL u \right) & = & u \circ \Pi_{\prec} \left( (a
  (u), a' (u)), \left( \mathcal{D} \DF \right) u \right) + \CC_T^{3 \alpha -
  2}\\
  & = & \Pi_{\precprec} (a (u), \vartheta) \circ \Pi_{\prec} \left( (a (u),
  a' (u)), \left( \mathcal{D} \DF \right) \vartheta \right) + \CC_T^{3 \alpha
  - 2}\\
  & = & \Pi_{\diamondsuit} \left( (a (u), a' (u)), \vartheta \circ \left(
  \mathcal{D} \DF \right) \vartheta \right) + \CC_T^{3 \alpha - 2}
\end{eqnarray*}


Finally the equation for $U^{\sharp}$ reads
\[ \LL U^{\sharp} = \Pi_{\diamondsuit} \left( (a (u), a' (u)), \vartheta \circ
   \mathcal{D} \Xi + \vartheta \circ \left( \mathcal{D} \DF \right) \vartheta
   \right) + \CC_T^{3 \alpha - 2} . \]
which can be solved essentially as we did in the simpler context. We see that
the general enhancement has the form
\[ \left( \xi, \vartheta \circ \mathcal{D} \Xi + \vartheta \circ \left(
   \mathcal{D} \DF \right) \vartheta \right) \]
which of course will require renormalization like we did before. In particular
\begin{eqnarray*}
  \left( \vartheta \circ \mathcal{D} \Xi + \vartheta \circ \left( \mathcal{D}
  \DF \right) \vartheta \right) (\eta, \eta') & = & \vartheta (\eta) \circ
  \eta'_2 \partial_{\eta_2} \xi (\eta_2, \cdot) + \vartheta (\eta) \circ
  \eta_1' \Delta \vartheta (\eta)\\
  & = & - \frac{\eta'_2}{\eta_1} (\Delta^{- 1} \xi (\eta_2, \cdot)) \circ
  \partial_{\eta_2} \xi (\eta_2, \cdot) + \frac{\eta_1'}{\eta_1^2} (\Delta^{-
  1} \xi (\eta_2, \cdot)) \circ \xi (\eta_2, \cdot)
\end{eqnarray*}
where we used that $\eta_1 \Delta \vartheta (\eta) = - \xi (\eta_2, \cdot)$.
Now observe that
\[ \mathbbm{E} [(\Delta^{- 1} \xi_{\varepsilon} (\eta_2, \cdot)) \circ
   \xi_{\varepsilon} (\eta_2, \cdot)] = - F (\eta_2, \eta_2)
   \sigma_{\varepsilon} \]
and that
\[ \mathbbm{E} [(\Delta^{- 1} \xi_{\varepsilon} (\eta_2, \cdot)) \circ
   \partial_{\eta_2} \xi_{\varepsilon} (\eta_2, \cdot)] = - (\partial_1 F)
   (\eta_2, \eta_2) \sigma_{\varepsilon} \]
with $\partial_1 F$ denoting the derivative with respect to the first entry.

In the end the renormalized enhanced noise is obtained as the limit in
$\mathcal{X}^{\alpha}$ of $(\xi_{\varepsilon}, \Xi_{2, \varepsilon})$ where
\[ \Xi_{2, \varepsilon} (\eta, \eta') = - \frac{\eta'_2}{\eta_1} (\Delta^{-
   1} \xi_{\varepsilon} (\eta_2, \cdot)) \circ \partial_{\eta_2}
   \xi_{\varepsilon} (\eta_2, \cdot) + \frac{\eta_1'}{\eta_1^2} (\Delta^{- 1}
   \xi_{\varepsilon} (\eta_2, \cdot)) \circ \xi_{\varepsilon} (\eta_2, \cdot)
   - H_{\varepsilon} (\eta, \eta') \]
with
\[ H_{\varepsilon} (\eta, \eta') = \frac{\eta'_2}{\eta_1} (\partial_1 F)
   (\eta_2, \eta_2) \sigma_{\varepsilon} - \frac{\eta_1'}{\eta_1^2} F (\eta_2,
   \eta_2) \sigma_{\varepsilon} . \]


We remark that if we take $F (\eta_2, \tilde{\eta}_2) = \eta_2 \tilde{\eta}_2$
we reobtain the situation treated in Section~\ref{sec:gen-nonlinear}, indeed
in this case
\[ \Pi_{\diamondsuit} ((a (u), a' (u)), H_{\varepsilon}) = \frac{a_2' (u) a_2
   (u)}{a_1 (u)} \sigma_{\varepsilon} - \frac{a'_1 (u) a_2 (u)^2}{a_1 (u)^2}
   \sigma_{\varepsilon} . \]
which coincides with~(\ref{eq:renorm-non-linear-rhs}).

\begin{remark}
  Consider the more general equation (\ref{eq:general-model}), where the noise
  depends explicitly on time, e.g. with a covariance
  \[ \mathbbm{E} [\xi (\eta, t, x) \xi (\eta', t', x')] = F (\eta, \eta') Q (t
     - t', x - x') \]
  with $F$ a smooth function and $Q$ a distribution of parabolic regularity
  $\rho > - 4 / 3$. First note that the coefficient $a_1 (u) \in [\lambda, 1]$
  in front of the time derivative can be eliminated trivially by dividing.
  
  In order to handle the time dependence of the noise, the framework of this
  paper will still apply, provided we consider space--time paraproducts
  instead of paraproducts which act only on the space variable. This can be
  done exactly following the lines of the
  paper~{\cite{gubinelli_paracontrolled_2015}} where time paraproducts were
  considered in the paracontrolled approach to solutions to SDE driven by
  gaussian signals.
  
  The constraint of regularity $\rho > - 4 / 3$ does allow to treat a noise
  which is white in time and smooth in space, but not a space--time white
  noise. It is well known that the first order paracontrolled approach on
  which the present paper is based does not allow to treat this kind of
  irregular signals in full generality.
\end{remark}

\appendix\section{Besov spaces}\label{appendix:besov}

In this Appendix we collect some classical results from harmonic analysis
needed in the paper. For a gentle introduction to Littlewood-Paley theory and
Besov spaces see the recent monograph~{\cite{bahouri_fourier_2011}}, where
most of our results are taken from. There the case of tempered distributions
on $\mathbbm{R}^d$ is considered. The Schauder estimates for the heat
semigroup are classical and can be found
in~{\cite{gubinelli_paracontrolled_2015,gubinelli_lectures_2015}}.

\

Fix $d \in \mathbbm{N}$ and denote by $\mathbbm{T}^d = (\mathbbm{R}/ (2 \pi
\mathbbm{Z}))^d$ the $d$-dimensional torus. We focus here on distributions and
SPDEs on the torus, but everything in this Appendix applies mutatis mutandis
on the full space $\mathbbm{R}^d$, see~{\cite{gubinelli_paracontrolled_2015}}.
The space of distributions $\CD' = \CD' (\mathbbm{T}^d)$ is defined as the set
of linear maps $f$ from $C^{\infty} = C^{\infty} (\mathbbm{T}^d, \mathbbm{C})$
to $\mathbbm{C}$, such that there exist $k \in \mathbbm{N}$ and $C > 0$ with
\[ | \langle f, \varphi \rangle | \assign | f (\varphi) | \leqslant C \sup_{|
   \mu | \leqslant k} \| \partial^{\mu} \varphi \|_{L^{\infty}
   (\mathbbm{T}^d)} \]
for all $\varphi \in C^{\infty}$. In particular, the Fourier transform $\CF f
: \mathbbm{Z}^d \rightarrow \mathbbm{C}$, $\CF f (k) = \langle f, e^{- i k
\cdummy} \rangle$, is defined for all $f \in \CD'$, and it satisfies $\left|
\CF f (k) \right| \leqslant | P (k) |$ for a suitable polynomial $P$. We will
also write $\hat{f} (k) = \CF f (k)$. Conversely, if $(g (k))_{k \in
\mathbbm{Z}^d}$ is at most of polynomial growth, then its inverse Fourier
transform
\[ \CF^{- 1} g = (2 \pi)^{- d} \sum_{k \in \mathbbm{Z}^d} e^{i \langle k,
   \cdummy \rangle} g (k) \]
defines a distribution, and we have $\CF^{- 1} \CF f = f$ as well as $\CF
\CF^{- 1} g = g$. To see this, it suffices to note that the Fourier transform
of $\varphi \in C^{\infty}$ decays faster than any rational function (we say
that it is of \tmtextit{rapid decay}). Indeed, for $\mu \in \mathbbm{N}^d_0$
we have $| k^{\mu} \hat{g} (k) | = \left| \CF (\partial^{\mu} g) (k) \right|
\leqslant \| \partial^{\mu} g \|_{L^1 (\mathbbm{T}^d)}$ for all $k \in
\mathbbm{Z}^d$. As a consequence we get the Parseval formula $\left\langle f,
\overline{\varphi} \right\rangle = (2 \pi)^{- d} \sum_k \hat{f} (k)
\overline{\hat{\varphi} (k)}$ for $f \in \CD'$ and $\varphi \in C^{\infty}$.

Linear maps on $\CD'$ can be defined by duality: if $A : C^{\infty}
\rightarrow C^{\infty}$ is such that for all $k \in \mathbbm{N}$ there exists
$n \in \mathbbm{N}$ and $C > 0$ with $\sup_{| \mu | \leqslant k} \|
\partial^{\mu} (A \varphi) \|_{L^{\infty}} \leqslant C \sup_{| \mu | \leqslant
n} \| \partial^{\mu} \varphi \|$, then we set $\langle^t A f, \varphi \rangle
= \langle f, A \varphi \rangle$. Differential operators are defined by
$\langle \partial^{\mu} f, \varphi \rangle = (- 1)^{| \mu |} \langle f,
\partial^{\mu} \varphi \rangle$. If $\varphi : \mathbbm{Z}^d \rightarrow
\mathbbm{C}$ grows at most polynomially, then it defines a \tmtextit{Fourier
multiplier}
\[ \varphi (\mathD) f = \CF^{- 1} \left( \varphi \CF f \right), \]
which gives us a distribution $\varphi (\mathD) f \in \CD'$ for every $f \in
\CD'$.

\

Littlewood-Paley blocks give a decomposition of any distribution on $\CD'$
into an infinite series of smooth functions.

\begin{definition}
  A \tmtextit{dyadic partition of unity} consists of two nonnegative radial
  functions $\chi, \rho \in C^{\infty} (\mathbbm{R}^d, \mathbbm{R})$, where
  $\rho$ is supported in a ball $\CB = \{ | x | \leqslant c \}$ and $\rho$ is
  supported in an annulus $\CA = \{ a \leqslant | x | \leqslant b \}$ for
  suitable $a, b, c > 0$, such that
  \begin{enumerate}
    \item $\chi + \sum_{j \geqslant 0} \rho (2^{- j} \cdummy) \equiv 1$ and
    
    \item \label{property-2-dyadic}$\chi \rho (2^{- j} \cdummy) \equiv 0$ for
    $j \geqslant 1$ and $\rho (2^{- i} \cdummy) \rho (2^{- j} \nosymbol
    \cdummy) \equiv 0$ for all $i, j \geqslant 0$ with $| i - j | > 1$.
  \end{enumerate}
  We will often write $\rho_{- 1} = \chi$ and $\rho_j = \rho (2^{- j}
  \cdummy)$ for $j \geqslant 0$.
\end{definition}

Dyadic partitions of unity exist, see~{\cite{bahouri_fourier_2011}}. The
reason for considering smooth partitions rather than indicator functions is
that indicator functions do not have good Fourier properties. We fix a dyadic
partition of unity $(\chi, \rho)$ and define the dyadic blocks
\[ \Delta_j f = \rho_j (\mathD) f = \CF^{- 1} (\rho_j \hat{f}), \quad j
   \geqslant - 1. \]
We also use the notation
\[ S_j f = \sum_{i < j - 1} \Delta_i f \]
and notice that
\[ \Delta_j f (x) = \int K_{i, x} (y) f (y) \mathd y, \qquad S_j f (x) = \int
   P_{i, x} (y) f (y) \mathd y \]
with $K_{i, x} (y) = 2^{di} K (2^i (x - y))$, $P_{i, x} (y) = \sum_{j < i - 1}
K_{j, x} (y)$, $K$ radial with zero mean.

Every dyadic block has a compactly supported Fourier transform and it belongs
therefore to~$C^{\infty}$. It is easy to see that $f = \sum_{j \geqslant - 1}
\Delta_j f = \lim_{j \rightarrow \infty} S_j f$ for all $f \in \CD'$.

For $\alpha \in \mathbbm{R}$, the H{\"o}lder-Besov space $\CC^{\alpha}$ is
given by $\CC^{\alpha} = B^{\alpha}_{\infty, \infty} (\mathbbm{T}^d,
\mathbbm{R})$, where for $p, q \in [1, \infty]$ we define
\[ B^{\alpha}_{p, q} = B^{\alpha}_{p, q} (\mathbbm{T}^d, \mathbbm{R}) =
   \left\{ f \in \DD' : \|f\|_{B^{\alpha}_{p, q}} = \left( \sum_{j \geqslant -
   1} (2^{j \alpha} \| \Delta_j f\|_{L^p})^q \right)^{1 / q} < \infty
   \right\}, \]
with the usual interpretation as $\ell^{\infty}$ norm in case $q = \infty$.
Then $B^{\alpha}_{p, q}$ is a Banach space and while the norm $\lVert \cdummy
\rVert_{B^{\alpha}_{p, q}}$ depends on $(\chi, \rho)$, the space
$B^{\alpha}_{p, q}$ does not, and any other dyadic partition of unity
corresponds to an equivalent norm.

If $\alpha \in (0, \infty) \setminus \mathbbm{N}$, then $\CC^{\alpha}$ is the
space of $\lfloor \alpha \rfloor$ times differentiable functions whose partial
derivatives of order $\lfloor \alpha \rfloor$ are ($\alpha - \lfloor \alpha
\rfloor$)-H{\"o}lder continuous (see page~99
of~{\cite{bahouri_fourier_2011}}). Note however, that for $k \in \mathbbm{N}$
the space $\CC^k$ is strictly larger than $C^k$, the space of $k$ times
continuously differentiable functions.

The following lemma gives useful characterisation of Besov regularity for
functions that can be decomposed into pieces which are localized in Fourier
space.

\begin{lemma}
  \label{lem: Besov regularity of series}{\tmdummy}
  
  \begin{enumeratenumeric}
    \item Let $\CA$ be an annulus, let $\alpha \in \mathbb{\mathbbm{R}}$, and
    let $(u_j)$ be a sequence of smooth functions such that $\CF u_j$ has its
    support in $2^j \CA$, and such that $\|u_j \|_{L^{\infty}} \lesssim 2^{- j
    \alpha}$ for all $j$. Then
    \[ u = \sum_{j \geqslant - 1} u_j \in \CC^{\alpha} \qquad \tmop{and}
       \qquad \|u\|_{\alpha} \lesssim \sup_{j \geqslant - 1} \{2^{j \alpha}
       \|u_j \|_{L^{\infty}} \} . \]
    \item Let $\CB$ be a ball, let $\alpha > 0$, and let $(u_j)$ be a sequence
    of smooth functions such that $\CF u_j$ has its support in $2^j \CB$, and
    such that $\|u_j \|_{L^{\infty}} \lesssim 2^{- j \alpha}$ for all $j$.
    Then
    \[ u = \sum_{j \geqslant - 1} u_j \in \CC^{\alpha} \qquad \tmop{and}
       \qquad \|u\|_{\alpha} \lesssim \sup_{j \geqslant - 1} \{2^{j \alpha}
       \|u_j \|_{L^{\infty}} \} . \]
  \end{enumeratenumeric}
\end{lemma}

\

The Bernstein inequalities of the next lemma are extremely useful when
dealing with functions with compactly supported Fourier transform.

\begin{lemma}
  \label{lemma:Bernstein}Let $\CA$ be an annulus and let $\CB$ be a ball. For
  any $k \in \mathbbm{N}_0$, $\lambda > 0$, and $1 \leqslant p \leqslant q
  \leqslant \infty$ we have that
  \begin{enumeratenumeric}
    \item if $u \in L^p$ is such that $\tmop{supp} (\CF u) \subseteq \lambda
    \CB$, then
    \[ \max_{\mu \in \mathbb{N}^d : | \mu | = k} \| \partial^{\mu} u\|_{L^q}
       \lesssim_k \lambda^{k + d \left( \frac{1}{p} - \frac{1}{q} \right)}
       \|u\|_{L^p} ; \]
    \item if $u \in L^p$ is such that $\tmop{supp} (\CF u) \subseteq \lambda
    \CA$, then
    \[ \lambda^k \|u\|_{L^p} \lesssim_k \max_{\mu \in \mathbb{N}^d : | \mu | =
       k} \| \partial^{\mu} u\|_{L^p} . \]
  \end{enumeratenumeric}
\end{lemma}

We recall the following standard heat kernel estimations.

\begin{lemma}[Schauder estimates]
  \label{l:schauder}Let $V_t = \int_0^t e^{\eta (t - s) \Delta} v_s \mathd s$
  and $\PP_t u_0 = e^{\eta \Delta t} u_0$, with $\eta \geqslant \lambda$. We
  define $\LL^{\alpha}_T$ and $C_{\eta}^k \LL^{\alpha}_T$, $C_{\eta}^k
  \CC^{\alpha}_T$ for $k \in \mathbbm{N}$ as in
  (\ref{e:def-Lnorm}),~(\ref{e:def-Cetanorm}) and introduce the norm
  \[ \| v \|_{\mathcal{E}^{\delta}_T \CC^{\alpha}} = \sup_{t \in [0, T]}
     t^{\delta} \| v (t, \cdummy) \|_{\CC^{\alpha}} . \]
  Then for any $\gamma \in [0, 1)$ and $\alpha \in \mathbbm{R}$:
  \begin{eqnarray*}
    \sup_{t \in [0, T]} t^{\gamma} \| V_t \|_{C_{\eta}^k \CC^{\alpha - 2
    \beta}} & \lesssim & \sup_{t \in [0, T]} t^{\gamma + \beta} \| v_t
    \|_{\CC^{\alpha - 2}} \quad \forall \beta \in [0, 1)\\
    \| V \|_{C_{\eta}^k \LL^{\alpha - 2 \beta}_T} & \lesssim & T^{\beta} \| v
    \|_{\CC^{\alpha - 2}_T} \quad \forall \beta \in [0, 1)\\
    \| V \|_{C_{\eta}^k \LL_T^{\alpha - 2 \beta}} & \lesssim & T^{\beta -
    \delta} \| v \|_{\mathcal{E}^{\delta}_{_T} \CC^{\alpha - 2}} \quad \forall
    \beta \in [0, 1), \forall \delta \in [0, \beta]\\
    \| V \|_{C_{\eta}^k \LL_T^{\gamma}} & \lesssim & T^{\frac{\rho -
    \gamma}{2} + 1 - \delta} \| v \|_{\mathcal{E}^{\delta}_{_T} \CC^{\rho}}
    \quad \forall \rho \in [\gamma - 2, \gamma), \forall \delta \in \left[ 0,
    \frac{\rho - \gamma}{2} + 1 \right]\\
    \left\| \PP_t u_0 \right\|_{C_{\eta}^k \CC^{\alpha}} & \lesssim & t^{-
    \frac{\alpha - \beta}{2}} \| u_0 \|_{\CC^{\beta}} \quad \forall \beta <
    \alpha\\
    \left\| \PP_t u_0 - \PP_s u_0 \right\|_{C_{\eta}^k \CC^{\alpha}} &
    \lesssim & s^{- \delta} | t - s |^{\frac{\beta - \alpha}{2} + \delta} \|
    u_0 \|_{\CC^{\beta}_{}}\\
    &  & \forall t \neq s \in \mathbbm{R}^+, \beta \leqslant \alpha + 2,
    \delta \in \left[ 0, 1 - \frac{\beta - \alpha}{2} \right]
  \end{eqnarray*}
\end{lemma}

We need the following interpolation lemma:

\begin{lemma}
  \label{l:interpol} Let $\gamma \in (0, 2)$, $0 < \varepsilon < \gamma$ and
  $u \in \LL^{\gamma}_T$ . Then
  \[ \| u \|_{C_T^{\gamma / 2 - \varepsilon / 2} L^{\infty}} \lesssim \| u
     \|_{\LL^{\gamma}_T} \]
\end{lemma}

\begin{proof}
  
  \begin{eqnarray*}
    \sup_{s \neq t} \frac{\| u_t - u_s \|_{L^{\infty}}}{| t - s |^{\gamma / 2
    - \varepsilon / 2}} & \leqslant & \sup_{s \neq t} \left[ \sum_{i \leqslant
    n} \frac{\| \Delta_i u_t - \Delta_i u_s \|_{L^{\infty}}}{| t - s |^{\gamma
    / 2 - \varepsilon / 2}} + \sum_{i > n} \frac{\| \Delta_i u_t - \Delta_i
    u_s \|_{L^{\infty}}}{| t - s |^{\gamma / 2 - \varepsilon / 2}} \right]
  \end{eqnarray*}
  and choosing $2^{- n - 1} \leqslant | t - s |^{1 / 2} \leqslant 2^{- n}$ we
  obtain
  \begin{eqnarray*}
    \sum_{i < n} \frac{\| \Delta_i u_t - \Delta_i u_s \|_{L^{\infty}}}{| t - s
    |^{\gamma / 2 - \varepsilon / 2}} & \lesssim & \| u \|_{C_T^{\gamma / 2}
    \CC^0} \sum_{i \leqslant n} | t - s |^{\varepsilon / 2}\\
    \sum_{i \geqslant n} \frac{\| \Delta_i u_t - \Delta_i u_s
    \|_{L^{\infty}}}{| t - s |^{\gamma / 2 - \varepsilon / 2}} & \lesssim & \|
    u \|_{\CC^{\gamma}_T} \sum_{i > n} 2^{- \gamma i} 2^{- (\gamma +
    \varepsilon) n}
  \end{eqnarray*}
  and this gives the result.
\end{proof}

Terms of the type $\| a (u (t, x)) \|_{\CC^0}$ with $a : \mathbbm{R}
\rightarrow \mathbbm{R}$ cannot be estimated directly with their H{\"o}lder
norm. In the following lemma we note some bounds used in
Section~\ref{s:fixpoint}.

\begin{lemma}
  Let $a \in C^3_b$ uniformly bounded and $u \in \LL^{\alpha}_T =
  \CC^{\alpha}_T \cap C_T^{\alpha / 2} \CC^0$, then
  \begin{eqnarray*}
    \| a (u) \|_{\LL^{\alpha}_T} & \lesssim & \| a \|_{L^{\infty}} + \| a'
    \|_{L^{\infty}} \left( \| u \|_{C_T^{\alpha / 2} \CC^0} + \| u
    \|_{\CC^{\alpha}_{_T}} \right) + \| a'' \|_{L^{\infty}} \| u
    \|_{C_T^{\alpha / 2} \CC^0} \| u \|_{\CC^{\alpha}_{_T}}\\
    & \lesssim & 1 + \| u \|_{\LL^{\alpha}_T} + \| u \|^2_{\LL_T^{\alpha}},
  \end{eqnarray*}
  \begin{eqnarray*}
    \| a (u_1) - a (u_2) \|_{\LL_T^{\alpha}} & \lesssim & \| u_1 - u_2
    \|_{\LL^{\alpha}_T} \left( 1 + \| u_1 \|_{\CC^{\alpha}_T} + \| u_2
    \|_{\CC^{\alpha}_T} \right)^2,\\
    \| a (u_1) - a (u_2) \|_{\CC_T^0} & \lesssim & \| u_1 - u_2 \|_{\CC^0_T}
    \left( 1 + \| u_1 \|_{\CC^{\alpha}_T} + \| u_2 \|_{\CC^{\alpha}_T} \right)
    .
  \end{eqnarray*}
\end{lemma}

\begin{proof}
  The bound on $\| a (u) \|_{\CC_T^{\alpha}}$ is trivial. We estimate $\| a
  (u_t) - a (u_s) \|_{\CC^0}$ as
  \[ \left| \int_z K_{i, x} (z) [a (u (t, z)) - a (u (s, z))] \right| = \left|
     \int_{z, \tau \in [0, 1]} K_{i, x} (z) a' (\delta_{\tau} u^{t z}_{s z})
     [u (t, z) - u (s, z)] \right| \]
  \begin{eqnarray*}
    & \leqslant & \left| \int_{z, w, \tau \in [0, 1]} K_{i, x} (z) \sum_{j
    \lesssim k \sim i} [\Delta_j u (t, z) - \Delta_j u (s, z)] K_{k, z} (w) a'
    (\delta_{\tau} u^{t w}_{s w}) \right|\\
    &  & + \left| \int_{z, w, \tau \in [0, 1]} K_{i, x} (z) \sum_{k \lesssim
    j \sim i} [\Delta_j u (t, z) - \Delta_j u (s, z)] K_{k, z} (w) a'
    (\delta_{\tau} u^{t w}_{s w}) \right|\\
    &  & + \left| \int_{z, w, \tau \in [0, 1]} K_{i, x} (z) \sum_{k \sim j
    \gtrsim i} [\Delta_j u (t, z) - \Delta_j u (s, z)] K_{k, z} (w) a'
    (\delta_{\tau} u^{t w}_{s w}) \right| .
  \end{eqnarray*}
  If $k > - 1$ we have
  \[ \int_w K_{k, z} (w) a' (\delta_{\tau} u^{t w}_{s w}) = \int_w K_{k, z}
     (w) [a' (\delta_{\tau} u^{t w}_{s w}) - a' (\delta_{\tau} u^{t z}_{s z})]
  \]
  and then the first term above becomes
  \[ \left| \int_{z, w, \tau, \sigma} K_{i, x} (z) \sum_{j \lesssim k \sim i}
     [\Delta_j u (t, z) - \Delta_j u (s, z)] K_{k, z} (w) a'' (\delta_{\sigma}
     (\delta_{\tau} u^t_s)^w_z) [\delta u^{t w}_{t z} - \delta u^{s w}_{s z}]
     \right| \]
  \[ \lesssim \| a'' \|_{L^{\infty}} \| u_t - u_s \|_{\CC^0} \sum_{j \lesssim
     k \sim i} \int_w | K_{k, z} (w) | | w - z |^{\alpha} \| u
     \|_{\CC^{\alpha}_T} \lesssim i 2^{- \alpha i} \| a'' \|_{L^{\infty}} \|
     u_t - u_s \|_{\CC^0} \| u \|_{\CC^{\alpha}_T} . \]
  The second term is
  \begin{eqnarray*}
    &  & \left| \int_{z, w, \tau \in [0, 1]} K_{i, x} (z) \sum_{j \sim i}
    [\Delta_j u (t, \cdummy) - \Delta_j u (s, \cdummy)] (z) P_{j, z} (w) a'
    (\delta_{\tau} u^{t w}_{s w}) \right|\\
    & \lesssim & \| u_t - u_s \|_{\CC^0} \| P_{j, z} (w) a' (\delta_{\tau}
    u^{t w}_{s w}) \|_{L^{\infty}} \lesssim \| u_t - u_s \|_{\CC^0} \| a'
    \|_{L^{\infty}}
  \end{eqnarray*}
  
  
  The third term can be estimated as the first one when $k > - 1$. Otherwise
  we just bound it as
  \[ \left| \int_{z, w, \tau \in [0, 1]} K_{i, x} (z) \sum_{j \leqslant N}
     [\Delta_j u (t, \cdummy) - \Delta_j u (s, \cdummy)] (z) K_{- 1, z} (w) a'
     (\delta_{\tau} u^{t w}_{s w}) \right| \lesssim \| u_t - u_s \|_{\CC^0} \|
     a' \|_{L^{\infty}} . \]
  
  
  For the three terms together we have the bound
  \[ \| a (u_t) - a (u_s) \|_{\CC^0} \lesssim \| u_t - u_s \|_{\CC^0} \left(
     \| a' \|_{L^{\infty}} + \| a'' \|_{L^{\infty}} \| u
     \|_{\CC^{\alpha}_{_T}} \right) \]
  With the same technique we obtain
  \begin{eqnarray*}
    \| a (u_1) - a (u_2) \|_{C_T^{\alpha / 2} \CC^0} & \lesssim & \| a'
    \|_{L^{\infty}} \| u_1 - u_2 \|_{C_T^{\alpha / 2} \CC^0}\\
    &  & + \| a'' \|_{L^{\infty}} \| u_1 - u_2 \|_{C_T^{\alpha / 2} \CC^0} \|
    u_1 - u_2 \|_{\CC^{\alpha}_T}\\
    &  & + \| a''' \|_{L^{\infty}} \| u_1 - u_2 \|_{C_T^{\alpha / 2} \CC^0}
    \| u_1 - u_2 \|^2_{\CC^{\alpha}_T}
  \end{eqnarray*}
  and this gives the second estimate. The third one can be obtained easily.
\end{proof}

\subsection{Bony's paraproduct}\label{ssec:bony}

\

In terms of Littlewood--Paley blocks, the product $fg$ can be decomposed as
\[ fg = \sum_{j \geqslant - 1} \sum_{i \geqslant - 1} \Delta_i f \Delta_j g =
   f \prec g + f \succ g + f \circ g. \]
Where the \tmtextit{paraproducts} $f \prec g$ and $f \succ g$ and the
\tmtextit{resonant product} $f \circ g$ are defined as
\[ f \prec g = g \succ f \assign \sum_{j \geqslant - 1} \sum_{i = - 1}^{j - 2}
   \Delta_i f \Delta_j g \qquad \text{and} \qquad f \circ g \assign \sum_{|i -
   j| \leqslant 1} \Delta_i f \Delta_j g. \]
We will often use the shortcuts $\sum_{i \sim j}$ for $\sum_{| i - j |
\leqslant 1}$ and $\sum_{i \lesssim j}$ for $\sum_{i < j - 1}$. Of course, the
decomposition depends on the dyadic partition of unity used to define the
blocks $\Delta_j$, and also on the particular choice of the pairs $(i, j)$ in
the diagonal part. The choice of taking all $(i, j)$ with $|i - j| \leqslant
1$ into the diagonal part corresponds to property~\ref{property-2-dyadic} in
our definition of dyadic partitions of unity.

Bony's crucial observation~{\cite{Bony1981,meyer_remarques_1981}} is that the
paraproduct $f \prec g$ (and thus $f \succ g$) is always a well-defined
distribution. Heuristically, $f \prec g$ behaves at large frequencies like $g$
(and thus retains the same regularity), and $f$ provides only a frequency
modulation of $g$. The only difficulty in constructing $fg$ for arbitrary
distributions lies in handling the resonant product $f \circ g$. The basic
result about these bilinear operations is given by the following estimates.

\begin{theorem}
  (Paraproduct estimates)\label{thm:paraproduct} For any $\beta \in
  \mathbbm{R}$ and $f, g \in \CD'$ we have
  \begin{equation}
    \label{eq:para-1} \|f \prec g\|_{\CC^{\beta}} \lesssim_{\beta}
    \|f\|_{L^{\infty}} \|g\|_{\CC^{\beta}},
  \end{equation}
  and for $\alpha < 0$ furthermore
  \begin{equation}
    \label{eq:para-2} \|f \prec g\|_{\CC^{\alpha + \beta}} \lesssim_{\alpha,
    \beta} \|f\|_{\CC^{\alpha}} \|g\|_{\CC^{\beta}} .
  \end{equation}
  For $\alpha + \beta > 0$ we have
  \begin{equation}
    \label{eq:para-3} \|f \circ g\|_{\CC^{\alpha + \beta}} \lesssim_{\alpha,
    \beta} \|f\|_{\CC^{\alpha}} \|g\|_{\CC^{\beta}} .
  \end{equation}
\end{theorem}

Bony proved also a basic \tmtextit{paralinearisation} result, soon after
improved by Meyer. We give here a particular version suited to our purposes.

\begin{theorem}
  \label{th:paralinearisation}Let $\alpha \in (0, 1)$, $f \in \left(
  \CC^{\alpha} \right)^k$ and $F \in C^3 (\mathbbm{R}^k ; \mathbbm{R})$ then
  \[ R_F (f) \assign F (f) - F' (f) \prec f \in \CC^{2 \alpha} \]
  with
  \[ \| R_F (f) \|_{\CC^{2 \alpha}} \lesssim \| F \|_{C^2} \left( 1 + \| f
     \|_{\CC^{\alpha}} \right)^2 . \]
  Moreover the map $f \mapsto R_F (f)$ is locally Lipshitz and
  \[ \| R_F (f) - R_F (\tilde{f}) \|_{\CC^{2 \alpha}} \lesssim \| F \|_{C^3}
     \left( 1 + \| f \|_{\CC^{\alpha}} + \| \tilde{f} \|_{\CC^{\alpha}}
     \right)^2 \| \tilde{f} - f \|_{\CC^{\alpha}} . \]
\end{theorem}

The additional key ingredient at the core of the paracontrolled approach is
the following commutation result proved
in~{\cite{gubinelli_paracontrolled_2015}}, Lemma~2.4:

\begin{lemma}
  \label{lem:commutator} Assume that $\alpha, \beta, \gamma \in \mathbb{R}$
  are such that $\alpha + \beta + \gamma > 0$ and $\beta + \gamma < 0$. Then
  for $f, g, h \in C^{\infty}$ the trilinear operator
  \[ C (f, g, h) \assign ((f \prec g) \circ h) - f (g \circ h) \]
  allows for the bound
  \begin{equation}
    \label{eq:commutator bound} \|C (f, g, h)\|_{\CC^{\beta + \gamma}}
    \lesssim \|f\|_{\CC^{\alpha}} \|g\|_{\CC^{\beta}} \|h\|_{\CC^{\gamma}},
  \end{equation}
  and can thus be uniquely extended to a bounded trilinear operator
  \[ C : \CC^{\alpha} \hspace{-0.17em} \times \hspace{-0.17em} \CC^{\beta}
     \hspace{-0.17em} \times \hspace{-0.17em} \CC^{\alpha} \rightarrow
     \CC^{\beta + \gamma} . \]
\end{lemma}

We will need the following two lemmas to compare standard and time-smoothed
paraproducts. The first one has essentially the same proof as
{\cite{gubinelli_paracontrolled_2015}}, Lemma~5.1.

\begin{lemma}
  \label{l:comm-paraproduct} Let $\rho \in (0, 2)$, $\gamma \in \mathbbm{R}$.
  Then for every $\varepsilon > 0$ we have the bound
  \[ \| g \prec h - g \precprec h \|_{\CC^{\rho + \gamma - \varepsilon}_T}
     \lesssim \| g \|_{C_T^{\rho / 2} \CC^0} \| h \|_{\CC^{\gamma}_T} . \]
\end{lemma}

The second lemma has a standard proof.

\begin{lemma}
  \label{l:comm-nonlin-para}Let $g \in \LL^{\rho}_T$, $h \in C_{\eta}^1
  \CC^{\gamma}_T$ with $\rho \in (0, 1)$, $\gamma \in \mathbbm{R}$. We have,
  $\forall \varepsilon > 0$
  \[ \| \Pi_{\prec} (g, h) - \Pi_{\precprec} (g, h) \|_{\CC^{\rho + \gamma -
     \varepsilon}_T} \lesssim \| g \|_{\LL^{\rho}_T} \| h \|_{C_{\eta}^1
     \CC^{\gamma}_T} . \]
\end{lemma}

\begin{thebibliography}{10}
  \bibitem[1]{alinhac_operateurs_1991}S.~Alinhac  and 
  P.~G{\'e}rard.{\newblock} \tmtextit{Op{\'e}rateurs pseudo-diff{\'e}rentiels
  et th{\'e}or{\`e}me de Nash-Moser}.{\newblock} EDP Sciences,
  1991.{\newblock}
  
  \bibitem[2]{allez_continuous_2015}R.~Allez  and  K.~Chouk.{\newblock} The
  continuous Anderson hamiltonian in dimension two.{\newblock}
  \tmtextit{ArXiv:1511.02718 [math]}, nov 2015.{\newblock} ArXiv:
  1511.02718.{\newblock}
  
  \bibitem[3]{bahouri_fourier_2011}H.~Bahouri, J.-Y.~Chemin , and 
  R.~Danchin.{\newblock} \tmtextit{Fourier Analysis and Nonlinear Partial
  Differential Equations}.{\newblock} Springer, jan 2011.{\newblock}
  
  \bibitem[4]{bailleul_heat_2015}I.~Bailleul  and  F.~Bernicot.{\newblock}
  Heat semigroup and singular PDEs.{\newblock} \tmtextit{ArXiv:1501.06822
  [math]}, jan 2015.{\newblock} ArXiv: 1501.06822.{\newblock}
  
  \bibitem[5]{bailleul_higher_2016}I.~Bailleul  and  F.~Bernicot.{\newblock}
  Higher order paracontrolled calculus.{\newblock} \tmtextit{ArXiv:1609.06966
  [math]}, sep 2016.{\newblock} ArXiv: 1609.06966.{\newblock}
  
  \bibitem[6]{bailleul_rough_2016}I.~Bailleul  and  R.~Catellier.{\newblock}
  Rough flows and homogenization in stochastic turbulence.{\newblock}
  \tmtextit{ArXiv:1601.01902 [math-ph]}, jan 2016.{\newblock} ArXiv:
  1601.01902.{\newblock}
  
  \bibitem[7]{hofmanova}Ismael Bailleul, Arnaud Debussche , and  Martina
  Hofmanova.{\newblock} Quasilinear generalized parabolic anderson model
  equation.{\newblock} \tmtextit{ArXiv:1610.06726 [math]}, Oct
  2016.{\newblock} ArXiv: 1610.06726.{\newblock}
  
  \bibitem[8]{bailleul_flows_2015}Isma{\"e}l Bailleul.{\newblock} Flows driven
  by rough paths.{\newblock} \tmtextit{Revista Matem{\'a}tica Iberoamericana},
  31(3):901--934, 2015.{\newblock}
  
  \bibitem[9]{Bony1981}J.-M.~Bony.{\newblock} Calcul symbolique et propagation
  des singularites pour les {\'e}quations aux d{\'e}riv{\'e}es partielles non
  lin{\'e}aires.{\newblock} \tmtextit{Ann. Sci. {\'E}c. Norm. Sup{\'e}r. (4)},
  14:209--246, 1981.{\newblock}
  
  \bibitem[10]{bony_microlocal_1991}J.-M.~Bony, G.~Grubb, L.~H{\"o}rmander,
  H.~Komatsu, J.~Sj{\"o}strand, L.~Cattabriga , and  L.~Rodino.{\newblock}
  \tmtextit{Microlocal Analysis and Applications Lectures given at the 2nd
  Session of the Centro Internazionale Matematico Estivo (C.I.M.E.) held at
  Montecatini Terme, Italy, July 3-11, 1989}.{\newblock} Springer Berlin
  Heidelberg, Berlin, Heidelberg, 1991.{\newblock}
  
  \bibitem[11]{catellier_paracontrolled_2013}R.~Catellier  and 
  K.~Chouk.{\newblock} Paracontrolled Distributions and the 3-dimensional
  Stochastic Quantization Equation.{\newblock} \tmtextit{ArXiv:1310.6869
  [math-ph]}, oct 2013.{\newblock} ArXiv: 1310.6869.{\newblock}
  
  \bibitem[12]{catellier_averaging_2016}R.~Catellier  and 
  M.~Gubinelli.{\newblock} Averaging along irregular curves and regularisation
  of ODEs.{\newblock} \tmtextit{Stochastic Processes and their Applications},
  126(8):2323--2366, aug 2016.{\newblock}
  
  \bibitem[13]{friz_course_2014}P.~K.~Friz  and  M.~Hairer.{\newblock}
  \tmtextit{A Course on Rough Paths: With an Introduction to Regularity
  Structures}.{\newblock} Springer, aug 2014.{\newblock}
  
  \bibitem[14]{gubinelli_rough_2012}M.~Gubinelli.{\newblock} Rough solutions
  for the periodic Korteweg--de Vries equation.{\newblock}
  \tmtextit{Communications on Pure and Applied Analysis}, 11(2):709--733,
  2012.{\newblock}
  
  \bibitem[15]{gubinelli_paracontrolled_2015}M.~Gubinelli, P.~Imkeller , and 
  N.~Perkowski.{\newblock} Paracontrolled distributions and singular
  PDEs.{\newblock} \tmtextit{Forum of Mathematics. Pi}, 3:0, 2015.{\newblock}
  
  \bibitem[16]{gubinelli_kpz_2015}M.~Gubinelli  and  N.~Perkowski.{\newblock}
  KPZ reloaded.{\newblock} \tmtextit{ArXiv:1508.03877 [math-ph]}, aug
  2015.{\newblock} ArXiv: 1508.03877.{\newblock}
  
  \bibitem[17]{gubinelli_lectures_2015}M.~Gubinelli  and 
  N.~Perkowski.{\newblock} Lectures on singular stochastic PDEs.{\newblock}
  \tmtextit{ArXiv:1502.00157 [math]}, jan 2015.{\newblock} ArXiv:
  1502.00157.{\newblock}
  
  \bibitem[18]{hairer_solving_2013}M.~Hairer.{\newblock} Solving the KPZ
  equation.{\newblock} \tmtextit{Annals of Mathematics}, 178(2):559--664, jun
  2013.{\newblock}
  
  \bibitem[19]{hairer_theory_2014}M.~Hairer.{\newblock} A theory of regularity
  structures.{\newblock} \tmtextit{Inventiones mathematicae}, 198(2):269--504,
  mar 2014.{\newblock}
  
  \bibitem[20]{hairer_singular_2014}M.~Hairer.{\newblock} Singular stochastic
  PDEs.{\newblock} \tmtextit{ArXiv:1403.6353 [math]}, mar 2014.{\newblock}
  ArXiv: 1403.6353.{\newblock}
  
  \bibitem[21]{hairer_class_2015}M.~Hairer  and  J.~Quastel.{\newblock} A
  class of growth models rescaling to KPZ.{\newblock}
  \tmtextit{ArXiv:1512.07845 [math-ph]}, dec 2015.{\newblock} ArXiv:
  1512.07845.{\newblock}
  
  \bibitem[22]{hairer_dynamical_2014}M.~Hairer  and  H.~Shen.{\newblock} The
  dynamical sine-Gordon model.{\newblock} \tmtextit{ArXiv:1409.5724
  [math-ph]}, sep 2014.{\newblock} ArXiv: 1409.5724.{\newblock}
  
  \bibitem[23]{hairer_large_2016}M.~Hairer  and  W.~Xu.{\newblock} Large scale
  behaviour of 3d phase coexistence models.{\newblock}
  \tmtextit{ArXiv:1601.05138 [math-ph]}, jan 2016.{\newblock} ArXiv:
  1601.05138.{\newblock}
  
  \bibitem[24]{hu_nonlinear_2014}Y.~Hu  and  K.~Le.{\newblock} Nonlinear Young
  integrals and differential systems in H\textbackslash "older
  media.{\newblock} \tmtextit{ArXiv:1404.7582 [math]}, apr 2014.{\newblock}
  
  \bibitem[25]{hormander_nash_moser_1990}L.~H{\"o}rmander.{\newblock} The
  Nash-Moser theorem and paradifferential operators.{\newblock} In
  \tmtextit{Analysis, et cetera},  pages  429--449. Academic Press, Boston,
  MA, 1990.{\newblock}
  
  \bibitem[26]{kelly_deterministic_2014}D.~Kelly  and 
  I.~Melbourne.{\newblock} Deterministic homogenization for fast-slow systems
  with chaotic noise.{\newblock} \tmtextit{ArXiv:1409.5748 [math]}, sep
  2014.{\newblock} ArXiv: 1409.5748.{\newblock}
  
  \bibitem[27]{kunita_stochastic_1984}H.~Kunita.{\newblock} Stochastic
  differential equations and stochastic flows of diffeomorphisms.{\newblock}
  In  P.~L.~Hennequin , editor , \tmtextit{{\'E}cole d'{\'E}t{\'e} de
  Probabilit{\'e}s de Saint-Flour XII - 1982},  number  1097  in  Lecture
  Notes in Mathematics,  pages  143--303. Springer Berlin Heidelberg,
  1984.{\newblock} DOI: 10.1007/BFb0099433.{\newblock}
  
  \bibitem[28]{lyons_differential_1998}T.~Lyons.{\newblock} Differential
  equations driven by rough signals.{\newblock} \tmtextit{Revista
  Matem{\'a}tica Iberoamericana},  pages  215--310, 1998.{\newblock}
  
  \bibitem[29]{lyons_differential_2007}T.~J.~Lyons, M.~J.~Caruana , and 
  T.~L{\'e}vy.{\newblock} \tmtextit{Differential Equations Driven by Rough
  Paths: Ecole d'Et{\'e} de Probabilit{\'e}s de Saint-Flour
  XXXIV-2004}.{\newblock} Springer, 1  edition , jun 2007.{\newblock}
  
  \bibitem[30]{lyons_system_2002}T.~Lyons  and  Z.~Qian.{\newblock}
  \tmtextit{System Control and Rough Paths}.{\newblock} Oxford University
  Press, 2002.{\newblock}
  
  \bibitem[31]{meyer_remarques_1981}Y.~Meyer.{\newblock} Remarques sur un
  th{\'e}or{\`e}me de J.-M. Bony.{\newblock} In \tmtextit{Rendiconti del
  Circolo Matematico di Palermo. Serie II},  pages  1--20. 1981.{\newblock}
  
  \bibitem[32]{mourrat_global_2016}J.-C.~Mourrat  and  H.~Weber.{\newblock}
  Global well-posedness of the dynamic \$\textbackslash
  Phi{\textasciicircum}4\_3\$ model on the torus.{\newblock}
  \tmtextit{ArXiv:1601.01234 [math]}, jan 2016.{\newblock} ArXiv:
  1601.01234.{\newblock}
  
  \bibitem[33]{otto_quasilinear_2016}F.~Otto  and  H.~Weber.{\newblock}
  Quasilinear SPDEs via rough paths.{\newblock} \tmtextit{ArXiv:1605.09744
  [math]}, may 2016.{\newblock} ArXiv: 1605.09744.{\newblock}
\end{thebibliography}

\end{document}
